%============================================================
%  Bd. 10 - Äquivalenzrelationen
%============================================================

\documentclass[main.tex]{subfiles}

\ifSubfilesClassLoaded{
  \BandSetupStandalone{B10}
}{
}

\title{Bd. 10 - Äquivalenzrelationen}
\author{Martin Kunze}
\date{}
\setcounter{file}{10}

\begin{document}

\title{Bd. 10 - Äquivalenzrelationen}
\author{Martin Kunze}
\date{}
\hypersetup{pageanchor=false}
\maketitle
\hypersetup{pageanchor=true}
\setcounter{tocdepth}{3}
\tableofcontents
\setcounter{file}{10}
\renewcommand{\FormulaBandID}{10}

\chapter{Einleitung}

Äquivalenzrelationen strukturieren einen Träger durch Klassen
ununterscheidbarer Elemente. Nach der allgemeinen Theorie folgen
Quotientenmengen, die zugehörigen Projektionsabbildungen und
Gleichmächtigkeit als wichtiges funktionsabhängiges Beispiel. Die
Ununterscheidbarkeitsrelation einer Mengenfamilie dient als durchgehendes
konkretes Beispiel.

\chapter{Äquivalenzrelationen und Quotienten}

\section{Begriff der Äquivalenzrelation}

\subsection{Axiome einer Äquivalenzrelation}

\FormulaAxiomDeltaK[Reflexivität]{%
  x\in A \vdash (x,x)\in R%
}{EqRelReflexivity}{
  \DeltaRow{Mengen}{A \dsep R \dsep x}
}

\FormulaAxiomDeltaK[Symmetrie]{%
  (x,y)\in R \vdash (y,x)\in R%
}{EqRelSymmetry}{
  \DeltaRow{Mengen}{A \dsep R \dsep x \dsep y}
}

\FormulaAxiomDeltaK[Transitivität]{%
  (x,y)\in R \dsep (y,z)\in R \vdash (x,z)\in R%
}{EqRelTransitivity}{
  \DeltaRow{Mengen}{A \dsep R \dsep x \dsep y \dsep z}
}

\subsection{Definition der Äquivalenzrelation}

\FormulaDefDeltaK[Begriff der Äquivalenzrelation]{\EqRel{R,A}}{EqRelDef}{
  \DeltaRow{Mengen}{A \dsep R \dsep x \dsep y \dsep z}
  \DeltaRow{\textbf{Axiome}}{}
  \DeltaRow{Relation auf \(A\)}
           {R \subseteq A \times A}
           [\FormulaRefAuto{F \subseteq A \times B}]
  \DeltaRow{Reflexivität}
           {x\in A \vdash (x,x)\in R}
           [\FormulaRefAuto{x\in A \vdash (x,x)\in R}]
  \DeltaRow{Symmetrie}
           {(x,y)\in R \vdash (y,x)\in R}
           [\FormulaRefAuto{(x,y)\in R \vdash (y,x)\in R}]
  \DeltaRow{Transitivität}
           {(x,y)\in R \dsep (y,z)\in R \vdash (x,z)\in R}
           [\FormulaRefAuto{(x,y)\in R \dsep (y,z)\in R \vdash (x,z)\in R}]
  \DeltaRow{\textbf{Neue Symbole}}{}
  \DeltaPrem{Äquivalenzrelationen}{}
}

\subsection{Beispiele von Äquivalenzrelationen}

\subsubsection{Ununterscheidbarkeitsrelation einer Mengenfamilie}

Ein wichtiges Beispiel entsteht aus einer Mengenfamilie \(M\). In der
Terminologie der Rough-Set-Theorie ist dies eine
Ununterscheidbarkeitsrelation (indiscernibility relation): Zwei Elemente des
Trägers \(\bigcup M\) sind bezüglich \(M\) ununterscheidbar, wenn sie in jedem
Mitglied von \(M\) entweder beide enthalten oder beide nicht enthalten sind.

\FormulaDefDeltaK[Ununterscheidbarkeitsrelation einer Mengenfamilie]{%
  \SameOccRel{M}\coloneqq
  \{(a,b)\in \bigcup M\times\bigcup M\mid
  \forall X\in M\,(a\in X\leftrightarrow b\in X)\}
}{SameOccRelDef}{%
  \DeltaRow{Mengen}{M\dsep X\dsep a\dsep b}
  \DeltaRow{Relationen}{\SameOccRel{M}}
  \DeltaRow{Neue Symbole}{\SameOccRel{M}}
}

\FormulaThmDeltaK[Charakterisierung der Ununterscheidbarkeitsrelation]{%
  (a,b)\in \SameOccRel{M}
  \eqvdash
  a\in\bigcup M\land b\in\bigcup M\land
  \forall X\in M\,(a\in X\leftrightarrow b\in X)
}{SameOccRelChar}{%
  \DeltaRow{Mengen}{M\dsep X\dsep a\dsep b}
  \DeltaRow{Relationen}{\SameOccRel{M}}
}
\begin{tabproofsplit}
\proofpart{\(\vdash\)}
  \proofstep{1}{(a,b)\in \SameOccRel{M}}{\rA}
  \proofstep{}{%
    \SameOccRel{M}\coloneqq
    \{(a,b)\in \bigcup M\times\bigcup M\mid
    \forall X\in M\,(a\in X\leftrightarrow b\in X)\}}{%
    \FormulaRefAuto{SameOccRelDef}}
  \proofstep{1}{(a,b)\in\bigcup M\times\bigcup M\land
    \forall X\in M\,(a\in X\leftrightarrow b\in X)}{%
    \FormulaRefAuto{x \in \{u \in A \mid P(u)\} \eqvdash x \in A \land P(x)}{2,1}}
  \proofstep{1}{(a,b)\in\bigcup M\times\bigcup M}{\rAEa{3}}
  \proofstep{1}{\forall X\in M\,(a\in X\leftrightarrow b\in X)}{\rAEb{3}}
  \proofstep{1}{a\in\bigcup M}{\FormulaRefAuto{(a,b)\in A\times B\vdash a\in A}{4}}
  \proofstep{1}{b\in\bigcup M}{\FormulaRefAuto{(a,b)\in A\times B\vdash b\in B}{4}}
  \proofstep{1}{a\in\bigcup M\land b\in\bigcup M\land
    \forall X\in M\,(a\in X\leftrightarrow b\in X)}{\rAI{6,\rAI{7,5}}}
\closeproofpart

\proofpart{\(\dashv\)}
  \proofstep{1}{a\in\bigcup M\land b\in\bigcup M\land
    \forall X\in M\,(a\in X\leftrightarrow b\in X)}{\rA}
  \proofstep{}{%
    \SameOccRel{M}\coloneqq
    \{(a,b)\in \bigcup M\times\bigcup M\mid
    \forall X\in M\,(a\in X\leftrightarrow b\in X)\}}{%
    \FormulaRefAuto{SameOccRelDef}}
  \proofstep{1}{a\in\bigcup M}{\rAEa{1}}
  \proofstep{1}{b\in\bigcup M}{\rAEa{\rAEb{1}}}
  \proofstep{1}{\forall X\in M\,(a\in X\leftrightarrow b\in X)}{\rAEb{\rAEb{1}}}
  \proofstep{1}{(a,b)\in\bigcup M\times\bigcup M}{%
    \FormulaRefAuto{a\in A\dsep b\in B\vdash (a,b)\in A\times B}{3,4}}
  \proofstep{1}{(a,b)\in\bigcup M\times\bigcup M\land
    \forall X\in M\,(a\in X\leftrightarrow b\in X)}{\rAI{6,5}}
  \proofstep{1}{(a,b)\in \SameOccRel{M}}{%
    \FormulaRefAuto{x \in \{u \in A \mid P(u)\} \eqvdash x \in A \land P(x)}{2,7}}
\closeproofpart
\end{tabproofsplit}

\FormulaThmDeltaK[Trägerbedingung der Ununterscheidbarkeitsrelation]{%
  \SameOccRel{M}\subseteq\bigcup M\times\bigcup M
}{SameOccRelCarrier}{%
  \DeltaRow{Mengen}{M\dsep X\dsep a\dsep b\dsep c}
  \DeltaRow{Relationen}{\SameOccRel{M}}
}
\begin{tabproofwide}
  \proofstepwidestar[1]{(a,b)\in \SameOccRel{M}}{\rA}
  \proofstepwidestar[1]{a\in\bigcup M\land b\in\bigcup M\land
    \forall X\in M\,(a\in X\leftrightarrow b\in X)}{%
    \FormulaRefAuto{SameOccRelChar}{1}}
  \proofstepwidestar[1]{a\in\bigcup M}{\rAEa{2}}
  \proofstepwidestar[1]{b\in\bigcup M}{\rAEa{\rAEb{2}}}
  \proofstepwidestar[1]{(a,b)\in\bigcup M\times\bigcup M}{%
    \FormulaRefAuto{a\in A\dsep b\in B\vdash (a,b)\in A\times B}{3,4}}
  \proofstepwidestar[]{\SameOccRel{M}\subseteq\bigcup M\times\bigcup M}{%
    \FormulaRefAuto{A\subseteq B \coloneq
    \forall x\,(x\in A \rightarrow x\in B)}{\rUI{\rRI{1,5}}}}
\end{tabproofwide}

\FormulaThmDeltaK[Reflexivität der Ununterscheidbarkeitsrelation]{%
  a\in\bigcup M \vdash (a,a)\in \SameOccRel{M}
}{SameOccRelReflexivity}{%
  \DeltaRow{Mengen}{M\dsep X\dsep a}
  \DeltaRow{Relationen}{\SameOccRel{M}}
}
\begin{tabproof}
  \proofstep{1}{a\in\bigcup M}{\rA}
  \proofstep{2}{X\in M}{\rA}
  \proofstep{1,2}{a\in X\leftrightarrow a\in X}{\FormulaRefAuto{P\leftrightarrow P}}
  \proofstep{1}{\forall X\in M\,(a\in X\leftrightarrow a\in X)}{\rUI{\rRI{2,3}}}
  \proofstep{1}{a\in\bigcup M\land a\in\bigcup M\land
    \forall X\in M\,(a\in X\leftrightarrow a\in X)}{\rAI{1,\rAI{1,4}}}
  \proofstep{1}{(a,a)\in \SameOccRel{M}}{%
    \FormulaRefAuto{SameOccRelChar}{5}}
\end{tabproof}

\FormulaThmDeltaK[Symmetrie der Ununterscheidbarkeitsrelation]{%
  (a,b)\in \SameOccRel{M} \vdash (b,a)\in \SameOccRel{M}
}{SameOccRelSymmetry}{%
  \DeltaRow{Mengen}{M\dsep X\dsep a\dsep b}
  \DeltaRow{Relationen}{\SameOccRel{M}}
}
\begin{tabproofwide}
  \proofstepwidestar[1]{(a,b)\in \SameOccRel{M}}{\rA}
  \proofstepwidestar[1]{a\in\bigcup M\land b\in\bigcup M\land
    \forall X\in M\,(a\in X\leftrightarrow b\in X)}{%
    \FormulaRefAuto{SameOccRelChar}{1}}
  \proofstepwidestar[1]{b\in\bigcup M}{\rAEa{\rAEb{2}}}
  \proofstepwidestar[1]{a\in\bigcup M}{\rAEa{2}}
  \proofstepwidestar[1]{\forall X\in M\,(b\in X\leftrightarrow a\in X)}{%
    \FormulaRefAuto{\forall x(P(x)\leftrightarrow Q(x))\vdash
    \forall x(Q(x)\leftrightarrow P(x))}{\rAEb{\rAEb{2}}}}
  \proofstepwidestar[1]{b\in\bigcup M\land a\in\bigcup M\land
    \forall X\in M\,(b\in X\leftrightarrow a\in X)}{\rAI{3,\rAI{4,5}}}
  \proofstepwidestar[1]{(b,a)\in \SameOccRel{M}}{%
    \FormulaRefAuto{SameOccRelChar}{6}}
\end{tabproofwide}

\FormulaThmDeltaK[Transitivität der Ununterscheidbarkeitsrelation]{%
  (a,b)\in \SameOccRel{M}\dsep (b,c)\in \SameOccRel{M}
  \vdash (a,c)\in \SameOccRel{M}
}{SameOccRelTransitivity}{%
  \DeltaRow{Mengen}{M\dsep X\dsep a\dsep b\dsep c}
  \DeltaRow{Relationen}{\SameOccRel{M}}
}
\begin{tabproofwide}
  \proofstepwidestar[1]{(a,b)\in \SameOccRel{M}}{\rA}
  \proofstepwidestar[2]{(b,c)\in \SameOccRel{M}}{\rA}
  \proofstepwidestar[1]{a\in\bigcup M\land b\in\bigcup M\land
    \forall X\in M\,(a\in X\leftrightarrow b\in X)}{%
    \FormulaRefAuto{SameOccRelChar}{1}}
  \proofstepwidestar[2]{b\in\bigcup M\land c\in\bigcup M\land
    \forall X\in M\,(b\in X\leftrightarrow c\in X)}{%
    \FormulaRefAuto{SameOccRelChar}{2}}
  \proofstepwidestar[1]{a\in\bigcup M}{\rAEa{3}}
  \proofstepwidestar[2]{c\in\bigcup M}{\rAEa{\rAEb{4}}}
  \proofstepwidestar[1,2]{\forall X\in M\,(a\in X\leftrightarrow c\in X)}{%
    \FormulaRefAuto{\forall x(P(x)\leftrightarrow Q(x)),
    \forall x(Q(x)\leftrightarrow R(x))\vdash
    \forall x(P(x)\leftrightarrow R(x))}{\rAEb{\rAEb{3}},\rAEb{\rAEb{4}}}}
  \proofstepwidestar[1,2]{a\in\bigcup M\land c\in\bigcup M\land
    \forall X\in M\,(a\in X\leftrightarrow c\in X)}{\rAI{5,\rAI{6,7}}}
  \proofstepwidestar[1,2]{(a,c)\in \SameOccRel{M}}{%
    \FormulaRefAuto{SameOccRelChar}{8}}
\end{tabproofwide}

\FormulaThmDeltaK[Ununterscheidbarkeitsrelation ist Äquivalenzrelation]{%
  \EqRel{\SameOccRel{M},\bigcup M}
}{SameOccRelEqRel}{%
  \DeltaRow{Mengen}{M\dsep X\dsep a\dsep b\dsep c}
  \DeltaRow{Relationen}{\SameOccRel{M}}
}
\begin{tabproof}
  \proofstep{}{\SameOccRel{M}\subseteq\bigcup M\times\bigcup M}{%
    \FormulaRefAuto{SameOccRelCarrier}}
  \proofstep{}{a\in\bigcup M \rightarrow (a,a)\in \SameOccRel{M}}{%
    \FormulaRefAuto{SameOccRelReflexivity}}
  \proofstep{}{(a,b)\in \SameOccRel{M} \rightarrow (b,a)\in \SameOccRel{M}}{%
    \FormulaRefAuto{SameOccRelSymmetry}}
  \proofstep{}{%
    \bigl((a,b)\in \SameOccRel{M}\land (b,c)\in \SameOccRel{M}\bigr)
    \rightarrow (a,c)\in \SameOccRel{M}}{%
    \FormulaRefAuto{SameOccRelTransitivity}}
  \proofstep{}{\EqRel{\SameOccRel{M},\bigcup M}}{%
    \FormulaRefAuto{EqRelDef}{1,2,3,4}}
\end{tabproof}

\subsection{Äquivalenzklassen}

\FormulaDefDeltaK[Äquivalenzklasse eines Elements]{%
  \EqRel{R,A}\dsep x\in A\vdash
  \EqClass{x}{R}{A}\coloneqq
  \{y\in A\mid (x,y)\in R\}
}{EqClassDef}{%
  \DeltaRow{Mengen}{A\dsep R\dsep x\dsep y}
  \DeltaRow{Relationen}{R}
  \DeltaRow{Neue Symbole}{\EqClass{x}{R}{A}}
}

\FormulaThmDeltaK[Charakterisierung der Äquivalenzklasse]{%
  \EqRel{R,A}\dsep x\in A
  \vdash
  y\in \EqClass{x}{R}{A}
  \leftrightarrow
  y\in A\land (x,y)\in R
}{EqClassChar}{%
  \DeltaRow{Mengen}{A\dsep R\dsep x\dsep y}
  \DeltaRow{Relationen}{R}
}
\begin{tabproof}
  \proofstep{1}{\EqRel{R,A}}{\rA}
  \proofstep{2}{x\in A}{\rA}
  \proofstep{1,2}{%
    \EqClass{x}{R}{A}\coloneqq
    \{y\in A\mid (x,y)\in R\}}{%
    \FormulaRefAuto{EqClassDef}{1,2}}

  \proofstep{4}{y\in \EqClass{x}{R}{A}}{\rA}
  \proofstep{1,2,4}{y\in A\land (x,y)\in R}{%
    \FormulaRefAuto{x \in \{u \in A \mid P(u)\} \eqvdash x \in A \land P(x)}{3,4}}
  \proofstep{1,2}{%
    y\in \EqClass{x}{R}{A}
    \rightarrow y\in A\land (x,y)\in R}{\rRI{4,5}}

  \proofstep{7}{y\in A\land (x,y)\in R}{\rA}
  \proofstep{1,2,7}{y\in \EqClass{x}{R}{A}}{%
    \FormulaRefAuto{x \in \{u \in A \mid P(u)\} \eqvdash x \in A \land P(x)}{3,7}}
  \proofstep{1,2}{%
    y\in A\land (x,y)\in R
    \rightarrow y\in \EqClass{x}{R}{A}}{\rRI{7,8}}

  \proofstep{1,2}{%
    y\in \EqClass{x}{R}{A}
    \leftrightarrow
    y\in A\land (x,y)\in R}{\rLRI{6,9}}
\end{tabproof}

\FormulaThmDeltaK[Repräsentant liegt in seiner Äquivalenzklasse]{%
  \EqRel{R,A}\dsep x\in A
  \vdash
  x\in \EqClass{x}{R}{A}
}{EqClassRepresentative}{%
  \DeltaRow{Mengen}{A\dsep R\dsep x}
  \DeltaRow{Relationen}{R}
}
\begin{tabproof}
  \proofstep{1}{\EqRel{R,A}}{\rA}
  \proofstep{2}{x\in A}{\rA}
  \proofstep{1,2}{(x,x)\in R}{\FormulaRefAuto{EqRelDef}{1,2}}
  \proofstep{1,2}{x\in A\land (x,x)\in R}{\rAI{2,3}}
  \proofstep{1,2}{x\in \EqClass{x}{R}{A}\leftrightarrow
    x\in A\land (x,x)\in R}{\FormulaRefAuto{EqClassChar}{1,2}}
  \proofstep{1,2}{x\in \EqClass{x}{R}{A}}{%
    \FormulaRefAuto{P \leftrightarrow Q, Q \vdash P}{5,4}}
\end{tabproof}

\FormulaThmDeltaK[Elemente einer Äquivalenzklasse liegen im Grundbereich]{%
  \EqRel{R,A}\dsep x\in A\dsep y\in \EqClass{x}{R}{A}
  \vdash y\in A
}{EqClassElementInCarrier}{%
  \DeltaRow{Mengen}{A\dsep R\dsep x\dsep y}
  \DeltaRow{Relationen}{R}
}
\begin{tabproof}
  \proofstep{1}{\EqRel{R,A}}{\rA}
  \proofstep{2}{x\in A}{\rA}
  \proofstep{3}{y\in \EqClass{x}{R}{A}}{\rA}
  \proofstep{1,2}{y\in \EqClass{x}{R}{A}\leftrightarrow
    y\in A\land (x,y)\in R}{\FormulaRefAuto{EqClassChar}{1,2}}
  \proofstep{1,2,3}{y\in A\land (x,y)\in R}{%
    \FormulaRefAuto{P \leftrightarrow Q, P \vdash Q}{4,3}}
  \proofstep{1,2,3}{y\in A}{\rAEa{5}}
\end{tabproof}

\FormulaThmDeltaK[Elemente einer Äquivalenzklasse sind äquivalent zum Repräsentanten]{%
  \EqRel{R,A}\dsep x\in A\dsep y\in \EqClass{x}{R}{A}
  \vdash (x,y)\in R
}{EqClassElementRelated}{%
  \DeltaRow{Mengen}{A\dsep R\dsep x\dsep y}
  \DeltaRow{Relationen}{R}
}
\begin{tabproof}
  \proofstep{1}{\EqRel{R,A}}{\rA}
  \proofstep{2}{x\in A}{\rA}
  \proofstep{3}{y\in \EqClass{x}{R}{A}}{\rA}
  \proofstep{1,2}{y\in \EqClass{x}{R}{A}\leftrightarrow
    y\in A\land (x,y)\in R}{\FormulaRefAuto{EqClassChar}{1,2}}
  \proofstep{1,2,3}{y\in A\land (x,y)\in R}{%
    \FormulaRefAuto{P \leftrightarrow Q, P \vdash Q}{4,3}}
  \proofstep{1,2,3}{(x,y)\in R}{\rAEb{5}}
\end{tabproof}

\FormulaThmDeltaK[Äquivalente Elemente haben dieselbe Äquivalenzklasse]{%
  \EqRel{R,A}\dsep x\in A\dsep y\in A\dsep (x,y)\in R
  \vdash
  \EqClass{x}{R}{A}=\EqClass{y}{R}{A}
}{EqClassEqualFromRel}{%
  \DeltaRow{Mengen}{A\dsep R\dsep x\dsep y\dsep z}
  \DeltaRow{Relationen}{R}
}
\begin{tabproofsplitwide}
  \proofpartwideindR[Teilmengenrichtung \(\EqClass{x}{R}{A}\subseteq\EqClass{y}{R}{A}\)]{%
    \EqRel{R,A}\dsep x\in A\dsep y\in A\dsep (x,y)\in R
    \vdash
    \EqClass{x}{R}{A}\subseteq\EqClass{y}{R}{A}
  }{%
    \EqRel{R,A}\dsep x\in A\dsep y\in A\dsep (x,y)\in R
    \vdash
    \EqClass{x}{R}{A}\subseteq\EqClass{y}{R}{A}
  }[EqClassEqualFromRelSubsetLeft]
    \proofstepwidestar[1]{\EqRel{R,A}}{\rA}
    \proofstepwidestar[2]{x\in A}{\rA}
    \proofstepwidestar[3]{y\in A}{\rA}
    \proofstepwidestar[4]{(x,y)\in R}{\rA}
    \proofstepwidestar[5]{z\in \EqClass{x}{R}{A}}{\rA}
    \proofstepwidestar[1,2,5]{z\in A}{%
      \FormulaRefAuto{EqClassElementInCarrier}{1,2,5}}
    \proofstepwidestar[1,2,5]{(x,z)\in R}{%
      \FormulaRefAuto{EqClassElementRelated}{1,2,5}}
    \proofstepwidestar[1,4]{(y,x)\in R}{\FormulaRefAuto{EqRelSymmetry}{4}}
    \proofstepwidestar[1,2,4,5]{(y,z)\in R}{%
      \FormulaRefAuto{(x,y)\in R \dsep (y,z)\in R \vdash (x,z)\in R}[ax]{8,7}}
    \proofstepwidestar[1,2,3,4,5]{z\in A\land (y,z)\in R}{\rAI{6,9}}
    \proofstepwidestar[1,3]{z\in \EqClass{y}{R}{A}\leftrightarrow
      z\in A\land (y,z)\in R}{\FormulaRefAuto{EqClassChar}{1,3}}
    \proofstepwidestar[1,2,3,4,5]{z\in \EqClass{y}{R}{A}}{%
      \FormulaRefAuto{P \leftrightarrow Q, Q \vdash P}{11,10}}
    \proofstepwidestar[1,2,3,4]{\EqClass{x}{R}{A}\subseteq\EqClass{y}{R}{A}}{%
      \FormulaRefAuto{A\subseteq B \coloneq
      \forall x\,(x\in A \rightarrow x\in B)}{\rUI{\rRI{5,12}}}}
  \closeproofpartwideind

  \proofpartwideindR[Teilmengenrichtung \(\EqClass{y}{R}{A}\subseteq\EqClass{x}{R}{A}\)]{%
    \EqRel{R,A}\dsep x\in A\dsep y\in A\dsep (x,y)\in R
    \vdash
    \EqClass{y}{R}{A}\subseteq\EqClass{x}{R}{A}
  }{%
    \EqRel{R,A}\dsep x\in A\dsep y\in A\dsep (x,y)\in R
    \vdash
    \EqClass{y}{R}{A}\subseteq\EqClass{x}{R}{A}
  }[EqClassEqualFromRelSubsetRight]
    \proofstepwidestar[1]{\EqRel{R,A}}{\rA}
    \proofstepwidestar[2]{x\in A}{\rA}
    \proofstepwidestar[3]{y\in A}{\rA}
    \proofstepwidestar[4]{(x,y)\in R}{\rA}
    \proofstepwidestar[5]{z\in \EqClass{y}{R}{A}}{\rA}
    \proofstepwidestar[1,3,5]{z\in A}{%
      \FormulaRefAuto{EqClassElementInCarrier}{1,3,5}}
    \proofstepwidestar[1,3,5]{(y,z)\in R}{%
      \FormulaRefAuto{EqClassElementRelated}{1,3,5}}
    \proofstepwidestar[1,3,4,5]{(x,z)\in R}{%
      \FormulaRefAuto{(x,y)\in R \dsep (y,z)\in R \vdash (x,z)\in R}[ax]{4,7}}
    \proofstepwidestar[1,2,3,4,5]{z\in A\land (x,z)\in R}{\rAI{6,8}}
    \proofstepwidestar[1,2]{z\in \EqClass{x}{R}{A}\leftrightarrow
      z\in A\land (x,z)\in R}{\FormulaRefAuto{EqClassChar}{1,2}}
    \proofstepwidestar[1,2,3,4,5]{z\in \EqClass{x}{R}{A}}{%
      \FormulaRefAuto{P \leftrightarrow Q, Q \vdash P}{10,9}}
    \proofstepwidestar[1,2,3,4]{\EqClass{y}{R}{A}\subseteq\EqClass{x}{R}{A}}{%
      \FormulaRefAuto{A\subseteq B \coloneq
      \forall x\,(x\in A \rightarrow x\in B)}{\rUI{\rRI{5,11}}}}
  \closeproofpartwideind

  \proofpartwide{Schluss}
    \proofstepwidestar[1]{\EqRel{R,A}}{\rA}
    \proofstepwidestar[2]{x\in A}{\rA}
    \proofstepwidestar[3]{y\in A}{\rA}
    \proofstepwidestar[4]{(x,y)\in R}{\rA}
    \proofstepwidestar[1,2,3,4]{\EqClass{x}{R}{A}\subseteq\EqClass{y}{R}{A}}{%
      \FormulaRefAuto{EqClassEqualFromRelSubsetLeft}{1,2,3,4}}
    \proofstepwidestar[1,2,3,4]{\EqClass{y}{R}{A}\subseteq\EqClass{x}{R}{A}}{%
      \FormulaRefAuto{EqClassEqualFromRelSubsetRight}{1,2,3,4}}
    \proofstepwidestar[1,2,3,4]{\EqClass{x}{R}{A}=\EqClass{y}{R}{A}}{%
      \FormulaRefAuto{A \subseteq B, B \subseteq A \vdash A = B}{5,6}}
  \closeproofpartwide
\end{tabproofsplitwide}

\FormulaThmDeltaK[Gleiche Äquivalenzklassen liefern Äquivalenz]{%
  \EqRel{R,A}\dsep x\in A\dsep y\in A\dsep
  \EqClass{x}{R}{A}=\EqClass{y}{R}{A}
  \vdash
  (x,y)\in R
}{EqClassRelFromEqual}{%
  \DeltaRow{Mengen}{A\dsep R\dsep x\dsep y}
  \DeltaRow{Relationen}{R}
}
\begin{tabproof}
  \proofstep{1}{\EqRel{R,A}}{\rA}
  \proofstep{2}{x\in A}{\rA}
  \proofstep{3}{y\in A}{\rA}
  \proofstep{4}{\EqClass{x}{R}{A}=\EqClass{y}{R}{A}}{\rA}
  \proofstep{1,2}{x\in \EqClass{x}{R}{A}}{%
    \FormulaRefAuto{EqClassRepresentative}{1,2}}
  \proofstep{1,2,4}{x\in \EqClass{y}{R}{A}}{\rIE{4,5}}
  \proofstep{1,2,3,4}{(y,x)\in R}{%
    \FormulaRefAuto{EqClassElementRelated}{1,3,6}}
  \proofstep{1,2,3,4}{(x,y)\in R}{\FormulaRefAuto{EqRelSymmetry}{7}}
\end{tabproof}

\FormulaThmDeltaK[Charakterisierung gleicher Äquivalenzklassen]{%
  \EqRel{R,A}\dsep x\in A\dsep y\in A
  \vdash
  \bigl(\EqClass{x}{R}{A}=\EqClass{y}{R}{A}\bigr)
  \leftrightarrow
  (x,y)\in R
}{EqClassEqualIffRel}{%
  \DeltaRow{Mengen}{A\dsep R\dsep x\dsep y}
  \DeltaRow{Relationen}{R}
}
\begin{tabproof}
  \proofstep{1}{\EqRel{R,A}}{\rA}
  \proofstep{2}{x\in A}{\rA}
  \proofstep{3}{y\in A}{\rA}

  \proofstep{4}{\EqClass{x}{R}{A}=\EqClass{y}{R}{A}}{\rA}
  \proofstep{1,2,3,4}{(x,y)\in R}{\FormulaRefAuto{EqClassRelFromEqual}{1,2,3,4}}
  \proofstep{1,2,3}{%
    \bigl(\EqClass{x}{R}{A}=\EqClass{y}{R}{A}\bigr)
    \rightarrow (x,y)\in R}{\rRI{4,5}}

  \proofstep{7}{(x,y)\in R}{\rA}
  \proofstep{1,2,3,7}{\EqClass{x}{R}{A}=\EqClass{y}{R}{A}}{%
    \FormulaRefAuto{EqClassEqualFromRel}{1,2,3,7}}
  \proofstep{1,2,3}{%
    (x,y)\in R
    \rightarrow \EqClass{x}{R}{A}=\EqClass{y}{R}{A}}{\rRI{7,8}}

  \proofstep{1,2,3}{%
    \bigl(\EqClass{x}{R}{A}=\EqClass{y}{R}{A}\bigr)
    \leftrightarrow (x,y)\in R}{\rLRI{6,9}}
\end{tabproof}

\subsection{Quotientenmengen}

\FormulaDefDeltaK[Quotientenmenge einer Äquivalenzrelation]{%
  \EqRel{R,A}\vdash
  \QuotientSet{A}{R}\coloneqq
  \{C\in\powerset(A)\mid
  \exists x\in A\,C=\EqClass{x}{R}{A}\}
}{QuotientSetDef}{%
  \DeltaRow{Mengen}{A\dsep R\dsep C\dsep x}
  \DeltaRow{Relationen}{R}
  \DeltaRow{Äquivalenzklassen}{\EqClass{x}{R}{A}}
  \DeltaRow{Neue Symbole}{\QuotientSet{A}{R}}
}

\FormulaThmDeltaK[Charakterisierung der Quotientenmenge]{%
  \EqRel{R,A}
  \vdash
  C\in \QuotientSet{A}{R}
  \leftrightarrow
  C\in\powerset(A)\land
  \exists x\in A\,C=\EqClass{x}{R}{A}
}{QuotientSetChar}{%
  \DeltaRow{Mengen}{A\dsep R\dsep C\dsep x}
  \DeltaRow{Relationen}{R}
  \DeltaRow{Äquivalenzklassen}{\EqClass{x}{R}{A}}
}
\begin{tabproof}
  \proofstep{1}{\EqRel{R,A}}{\rA}
  \proofstep{1}{%
    \QuotientSet{A}{R}\coloneqq
    \{C\in\powerset(A)\mid
    \exists x\in A\,C=\EqClass{x}{R}{A}\}}{%
    \FormulaRefAuto{QuotientSetDef}{1}}

  \proofstep{3}{C\in \QuotientSet{A}{R}}{\rA}
  \proofstep{1,3}{C\in\powerset(A)\land
    \exists x\in A\,C=\EqClass{x}{R}{A}}{%
    \FormulaRefAuto{x \in \{u \in A \mid P(u)\} \eqvdash x \in A \land P(x)}{2,3}}
  \proofstep{1}{C\in \QuotientSet{A}{R}
    \rightarrow C\in\powerset(A)\land
    \exists x\in A\,C=\EqClass{x}{R}{A}}{\rRI{3,4}}

  \proofstep{6}{C\in\powerset(A)\land
    \exists x\in A\,C=\EqClass{x}{R}{A}}{\rA}
  \proofstep{1,6}{C\in \QuotientSet{A}{R}}{%
    \FormulaRefAuto{x \in \{u \in A \mid P(u)\} \eqvdash x \in A \land P(x)}{2,6}}
  \proofstep{1}{C\in\powerset(A)\land
    \exists x\in A\,C=\EqClass{x}{R}{A}
    \rightarrow C\in \QuotientSet{A}{R}}{\rRI{6,7}}

  \proofstep{1}{C\in \QuotientSet{A}{R}
    \leftrightarrow
    C\in\powerset(A)\land
    \exists x\in A\,C=\EqClass{x}{R}{A}}{\rLRI{5,8}}
\end{tabproof}

\FormulaThmDeltaK[Äquivalenzklasse liegt in der Quotientenmenge]{%
  \EqRel{R,A}\dsep x\in A
  \vdash
  \EqClass{x}{R}{A}\in \QuotientSet{A}{R}
}{EqClassInQuotient}{%
  \DeltaRow{Mengen}{A\dsep R\dsep x\dsep y}
  \DeltaRow{Relationen}{R}
  \DeltaRow{Äquivalenzklassen}{\EqClass{x}{R}{A}}
}
\begin{tabproof}
  \proofstep{1}{\EqRel{R,A}}{\rA}
  \proofstep{2}{x\in A}{\rA}
  \proofstep{3}{y\in \EqClass{x}{R}{A}}{\rA}
  \proofstep{1,2,3}{y\in A}{%
    \FormulaRefAuto{EqClassElementInCarrier}{1,2,3}}
  \proofstep{1,2}{\EqClass{x}{R}{A}\subseteq A}{%
    \FormulaRefAuto{A\subseteq B \coloneq
    \forall x\,(x\in A \rightarrow x\in B)}{\rUI{\rRI{3,4}}}}
  \proofstep{1,2}{\EqClass{x}{R}{A}\in\powerset(A)}{%
    \FormulaRefAuto{x \in \powerset(A)\;\eqvdash\; x \subseteq A}{5}}
  \proofstep{}{ \EqClass{x}{R}{A}=\EqClass{x}{R}{A}}{\rII}
  \proofstep{2}{\exists z\in A\,\EqClass{x}{R}{A}=\EqClass{z}{R}{A}}{\rEI{\rAI{2,7}}}
  \proofstep{1,2}{%
    \EqClass{x}{R}{A}\in\powerset(A)\land
    \exists z\in A\,\EqClass{x}{R}{A}=\EqClass{z}{R}{A}}{\rAI{6,8}}
  \proofstep{1,2}{\EqClass{x}{R}{A}\in \QuotientSet{A}{R}}{%
    \FormulaRefAuto{QuotientSetChar}{1,9}}
\end{tabproof}

\FormulaThmDelta[Elemente der Quotientenmenge sind Teilmengen]{%
  \EqRel{R,A}\dsep C\in \QuotientSet{A}{R}
  \vdash
  C\subseteq A
}{%
  \DeltaRow{Mengen}{A\dsep R\dsep C}
  \DeltaRow{Relationen}{R}
}
\begin{tabproof}
  \proofstep{1}{\EqRel{R,A}}{\rA}
  \proofstep{2}{C\in \QuotientSet{A}{R}}{\rA}
  \proofstep{1,2}{%
    C\in\powerset(A)\land
    \exists x\in A\,C=\EqClass{x}{R}{A}}{%
    \FormulaRefAuto{QuotientSetChar}{1,2}}
  \proofstep{1,2}{C\in\powerset(A)}{\rAEa{3}}
  \proofstep{1,2}{C\subseteq A}{%
    \FormulaRefAuto{x \in \powerset(A)\;\eqvdash\; x \subseteq A}{4}}
\end{tabproof}

\FormulaThmDeltaK[Elemente der Quotientenmenge haben einen Repräsentanten]{%
  \EqRel{R,A}\dsep C\in \QuotientSet{A}{R}
  \vdash
  \exists x\in A\,C=\EqClass{x}{R}{A}
}{QuotientRepresentative}{%
  \DeltaRow{Mengen}{A\dsep R\dsep C\dsep x}
  \DeltaRow{Relationen}{R}
}
\begin{tabproof}
  \proofstep{1}{\EqRel{R,A}}{\rA}
  \proofstep{2}{C\in \QuotientSet{A}{R}}{\rA}
  \proofstep{1,2}{%
    C\in\powerset(A)\land
    \exists x\in A\,C=\EqClass{x}{R}{A}}{%
    \FormulaRefAuto{QuotientSetChar}{1,2}}
  \proofstep{1,2}{\exists x\in A\,C=\EqClass{x}{R}{A}}{\rAEb{3}}
\end{tabproof}


\section{Quotientenprojektionen als Funktionen}

\subsection{Quotientenprojektionen und Quotientenbildabbildungen}

\FormulaDefDeltaK[Funktionsvorschrift der Quotientenprojektion]{%
  \EqRel{R,A}\dsep x\in A\vdash
  \tQuotProj{A}{R}(x)\coloneqq \EqClass{x}{R}{A}
}{QuotProjRuleDef}{%
  \DeltaRow{Mengen}{A\dsep R\dsep x}
  \DeltaRow{Relationen}{R}
  \DeltaRow{Äquivalenzklassen}{\EqClass{x}{R}{A}}
  \DeltaRow{Neue Symbole}{\tQuotProj{A}{R}}
}

\FormulaThmDeltaK[Typisierung der Quotientenprojektionsvorschrift]{%
  \EqRel{R,A}\dsep x\in A\vdash
  \tQuotProj{A}{R}(x)\in\QuotientSet{A}{R}
}{QuotProjRuleInQuotient}{%
  \DeltaRow{Mengen}{A\dsep R\dsep x}
  \DeltaRow{Relationen}{R}
  \DeltaRow{Quotientenmenge}{\QuotientSet{A}{R}}
}
\begin{tabproof}
  \proofstep{1}{\EqRel{R,A}}{\rA}
  \proofstep{2}{x\in A}{\rA}
  \proofstep{1,2}{\tQuotProj{A}{R}(x)=\EqClass{x}{R}{A}}{%
    \FormulaRefAuto{QuotProjRuleDef}{1,2}}
  \proofstep{1,2}{\EqClass{x}{R}{A}\in\QuotientSet{A}{R}}{%
    \FormulaRefAuto{EqClassInQuotient}{1,2}}
  \proofstep{1,2}{\tQuotProj{A}{R}(x)\in\QuotientSet{A}{R}}{\rIE{3,4}}
\end{tabproof}

\FormulaThmDeltaK[Existenz der Quotientenprojektion]{%
  \EqRel{R,A}\vdash
  \exists! q\colon A\to\QuotientSet{A}{R}\,
  \forall x\in A\,q(x)=\tQuotProj{A}{R}(x)
}{QuotProjExists}{%
  \DeltaRow{Mengen}{A\dsep R\dsep x}
  \DeltaRow{Funktionen}{q}
  \DeltaRow{Funktionensymbol}{\tQuotProj{A}{R}}
}
\begin{tabproof}
  \proofstep{1}{\EqRel{R,A}}{\rA}
  \proofstep{1}{%
    \begin{aligned}[t]
      &\exists! q\colon A\to\QuotientSet{A}{R}\,{}\\
      &\forall x\in A\,q(x)=\tQuotProj{A}{R}(x)
    \end{aligned}}{%
    \FormulaRefAuto{%
      x\in A\vdash t(x)\coloneq y\dsep x\in A\vdash t(x)\in B
      \exists! F\colon A\to B\,\forall x\in A\,F(x)=t(x)
    }{%
      \FormulaRefAuto{QuotProjRuleDef}{1},
      \FormulaRefAuto{QuotProjRuleInQuotient}{1}}}
\end{tabproof}

\FormulaDefDeltaK[Quotientenprojektion]{%
  \EqRel{R,A}\vdash
  \QuotProj{A}{R}\coloneqq
  \iota q\Bigl(q\colon A\to\QuotientSet{A}{R}\land
  \forall x\in A\,q(x)=\tQuotProj{A}{R}(x)\Bigr)
}{QuotProjDef}{%
  \DeltaRow{Mengen}{A\dsep R\dsep x}
  \DeltaRow{Funktionen}{q\dsep \QuotProj{A}{R}}
  \DeltaRow{Funktionensymbol}{\tQuotProj{A}{R}}
  \DeltaRow{Neue Symbole}{\QuotProj{A}{R}}
}

\FormulaThmDeltaK[Typisierung der Quotientenprojektion]{%
  \EqRel{R,A}\vdash
  \QuotProj{A}{R}\colon A\to\QuotientSet{A}{R}
}{QuotProjFunction}{%
  \DeltaRow{Mengen}{A\dsep R}
  \DeltaRow{Funktionen}{\QuotProj{A}{R}}
}
\begin{tabproof}
  \proofstep{1}{\EqRel{R,A}}{\rA}
  \proofstep{1}{\QuotProj{A}{R}\colon A\to\QuotientSet{A}{R}}{%
    \rAEa{\FormulaRefAuto{QuotProjDef}{1}}}
\end{tabproof}

\FormulaThmDeltaK[Wert der Quotientenprojektion]{%
  \EqRel{R,A}\dsep x\in A\vdash
  \QuotProj{A}{R}(x)=\EqClass{x}{R}{A}
}{QuotProjValue}{%
  \DeltaRow{Mengen}{A\dsep R\dsep x}
  \DeltaRow{Funktionen}{\QuotProj{A}{R}}
  \DeltaRow{Äquivalenzklassen}{\EqClass{x}{R}{A}}
}
\begin{tabproof}
  \proofstep{1}{\EqRel{R,A}}{\rA}
  \proofstep{2}{x\in A}{\rA}
  \proofstep{1}{\forall x\in A\,
    \QuotProj{A}{R}(x)=\tQuotProj{A}{R}(x)}{%
    \rAEb{\FormulaRefAuto{QuotProjDef}{1}}}
  \proofstep{1,2}{\QuotProj{A}{R}(x)=\tQuotProj{A}{R}(x)}{%
    \rRE{2,\rUE{3}}}
  \proofstep{1,2}{\tQuotProj{A}{R}(x)=\EqClass{x}{R}{A}}{%
    \FormulaRefAuto{QuotProjRuleDef}{1,2}}
  \proofstep{1,2}{\QuotProj{A}{R}(x)=\EqClass{x}{R}{A}}{%
    \FormulaRefAuto{a = b,\, b = c \vdash a = c}{4,5}}
\end{tabproof}

\FormulaThmDeltaK[Surjektivität der Quotientenprojektion]{%
  \EqRel{R,A}\vdash
  \QuotProj{A}{R}\colon A\sur\QuotientSet{A}{R}
}{QuotProjSurjective}{%
  \DeltaRow{Mengen}{A\dsep R\dsep C\dsep x}
  \DeltaRow{Relationen}{R}
  \DeltaRow{Funktionen}{\QuotProj{A}{R}}
}
\begin{tabproof}
  \proofstep{1}{\EqRel{R,A}}{\rA}
  \proofstep{1}{\QuotProj{A}{R}\colon A\to\QuotientSet{A}{R}}{%
    \FormulaRefAuto{QuotProjFunction}{1}}
  \proofstep{3}{C\in\QuotientSet{A}{R}}{\rA}
  \proofstep{1,3}{\exists x\in A\,C=\EqClass{x}{R}{A}}{%
    \FormulaRefAuto{QuotientRepresentative}{1,3}}
  \proofstep{5}{x\in A\land C=\EqClass{x}{R}{A}}{\rA}
  \proofstep{5}{x\in A}{\rAEa{5}}
  \proofstep{5}{C=\EqClass{x}{R}{A}}{\rAEb{5}}
  \proofstep{1,5}{\QuotProj{A}{R}(x)=\EqClass{x}{R}{A}}{%
    \FormulaRefAuto{QuotProjValue}{1,6}}
  \proofstep{1,5}{C=\QuotProj{A}{R}(x)}{%
    \FormulaRefAuto{a = b,\, c = b \vdash a = c}{7,8}}
  \proofstep{1,5}{\QuotProj{A}{R}(x)=C}{%
    \FormulaRefAuto{a=b\vdash b=a}{9}}
  \proofstep{1,5}{\exists x\in A\,\QuotProj{A}{R}(x)=C}{%
    \rEI{\rAI{6,10}}}
  \proofstep{1,3}{\exists x\in A\,\QuotProj{A}{R}(x)=C}{%
    \rEE{4,5,11}}
  \proofstep{1}{\forall C\in\QuotientSet{A}{R}\,
    \exists x\in A\,\QuotProj{A}{R}(x)=C}{%
    \rUI{\rRI{3,12}}}
  \proofstep{1}{\QuotProj{A}{R}\colon A\sur\QuotientSet{A}{R}}{%
    \FormulaRefAuto{Surjektive Funktion}{2,13}}
\end{tabproof}

\FormulaThmDeltaK[Gleichheit von Projektionswerten]{%
  \EqRel{R,A}\dsep x\in A\dsep y\in A
  \vdash
  \bigl(\QuotProj{A}{R}(x)=\QuotProj{A}{R}(y)\bigr)
  \leftrightarrow (x,y)\in R
}{QuotProjEqualIffRel}{%
  \DeltaRow{Mengen}{A\dsep R\dsep x\dsep y}
  \DeltaRow{Relationen}{R}
  \DeltaRow{Funktionen}{\QuotProj{A}{R}}
}
\begin{tabproof}
  \proofstep{1}{\EqRel{R,A}}{\rA}
  \proofstep{2}{x\in A}{\rA}
  \proofstep{3}{y\in A}{\rA}
  \proofstep{1,2}{\QuotProj{A}{R}(x)=\EqClass{x}{R}{A}}{%
    \FormulaRefAuto{QuotProjValue}{1,2}}
  \proofstep{1,3}{\QuotProj{A}{R}(y)=\EqClass{y}{R}{A}}{%
    \FormulaRefAuto{QuotProjValue}{1,3}}

  \proofstep{6}{\QuotProj{A}{R}(x)=\QuotProj{A}{R}(y)}{\rA}
  \proofstep{1,2}{\EqClass{x}{R}{A}=\QuotProj{A}{R}(x)}{%
    \FormulaRefAuto{a=b\vdash b=a}{4}}
  \proofstep{1,2,3,6}{\EqClass{x}{R}{A}=\QuotProj{A}{R}(y)}{%
    \FormulaRefAuto{a = b,\, b = c \vdash a = c}{7,6}}
  \proofstep{1,2,3,6}{\EqClass{x}{R}{A}=\EqClass{y}{R}{A}}{%
    \FormulaRefAuto{a = b,\, b = c \vdash a = c}{8,5}}
  \proofstep{1,2,3,6}{(x,y)\in R}{%
    \FormulaRefAuto{EqClassRelFromEqual}{1,2,3,9}}
  \proofstep{1,2,3}{%
    \QuotProj{A}{R}(x)=\QuotProj{A}{R}(y)\rightarrow (x,y)\in R}{%
    \rRI{6,10}}

  \proofstep{12}{(x,y)\in R}{\rA}
  \proofstep{1,2,3,12}{\EqClass{x}{R}{A}=\EqClass{y}{R}{A}}{%
    \FormulaRefAuto{EqClassEqualFromRel}{1,2,3,12}}
  \proofstep{1,2,3,12}{\QuotProj{A}{R}(x)=\EqClass{y}{R}{A}}{%
    \FormulaRefAuto{a = b,\, b = c \vdash a = c}{4,13}}
  \proofstep{1,3}{\EqClass{y}{R}{A}=\QuotProj{A}{R}(y)}{%
    \FormulaRefAuto{a=b\vdash b=a}{5}}
  \proofstep{1,2,3,12}{\QuotProj{A}{R}(x)=\QuotProj{A}{R}(y)}{%
    \FormulaRefAuto{a = b,\, b = c \vdash a = c}{14,15}}
  \proofstep{1,2,3}{%
    (x,y)\in R\rightarrow
    \QuotProj{A}{R}(x)=\QuotProj{A}{R}(y)}{%
    \rRI{12,16}}

  \proofstep{1,2,3}{%
    \bigl(\QuotProj{A}{R}(x)=\QuotProj{A}{R}(y)\bigr)
    \leftrightarrow (x,y)\in R}{\rLRI{11,17}}
\end{tabproof}

\FormulaDefDeltaK[Quotientenbildabbildung einer Teilfamilie]{%
  \EqRel{R,A}\dsep M\subseteq\powerset(A)\vdash
  \QuotFamMap{A}{R}{M}\coloneqq
  \PowMap{\QuotProj{A}{R}}\restriction_M
}{QuotFamMapDef}{%
  \DeltaRow{Mengen}{A\dsep R\dsep M}
  \DeltaRow{Funktionen}{\QuotProj{A}{R}\dsep \PowMap{\QuotProj{A}{R}}\dsep
  \QuotFamMap{A}{R}{M}}
  \DeltaRow{Neue Symbole}{\QuotFamMap{A}{R}{M}}
}

\FormulaThmDeltaK[Typisierung der Quotientenbildabbildung]{%
  \EqRel{R,A}\dsep M\subseteq\powerset(A)\vdash
  \QuotFamMap{A}{R}{M}\colon
  M\to\powerset(\QuotientSet{A}{R})
}{QuotFamMapFunction}{%
  \DeltaRow{Mengen}{A\dsep R\dsep M}
  \DeltaRow{Funktionen}{\QuotProj{A}{R}\dsep \QuotFamMap{A}{R}{M}}
}
\begin{tabproof}
  \proofstep{1}{\EqRel{R,A}}{\rA}
  \proofstep{2}{M\subseteq\powerset(A)}{\rA}
  \proofstep{1}{\QuotProj{A}{R}\colon A\to\QuotientSet{A}{R}}{%
    \FormulaRefAuto{QuotProjFunction}{1}}
  \proofstep{1}{\PowMap{\QuotProj{A}{R}}\colon
    \powerset(A)\to\powerset(\QuotientSet{A}{R})}{%
    \FormulaRefAuto{F\colon A\to B \vdash
    \PowMap{F}\colon \powerset(A)\to \powerset(B)}{3}}
  \proofstep{1,2}{\PowMap{\QuotProj{A}{R}}\restriction_M
    \colon M\to\powerset(\QuotientSet{A}{R})}{%
    \FormulaRefAuto{C\subseteq A \vdash F\restriction_C\colon C\to B}{2,4}}
  \proofstep{1,2}{\QuotFamMap{A}{R}{M}=
    \PowMap{\QuotProj{A}{R}}\restriction_M}{%
    \FormulaRefAuto{QuotFamMapDef}{1,2}}
  \proofstep{1,2}{\QuotFamMap{A}{R}{M}\colon
    M\to\powerset(\QuotientSet{A}{R})}{\rIE{6,5}}
\end{tabproof}

\FormulaThmDeltaK[Wert der Quotientenbildabbildung]{%
  \EqRel{R,A}\dsep M\subseteq\powerset(A)\dsep X\in M\vdash
  \QuotFamMap{A}{R}{M}(X)=\QuotProj{A}{R}[X]
}{QuotFamMapValue}{%
  \DeltaRow{Mengen}{A\dsep R\dsep M\dsep X}
  \DeltaRow{Funktionen}{\QuotProj{A}{R}\dsep \QuotFamMap{A}{R}{M}}
}
\begin{tabproof}
  \proofstep{1}{\EqRel{R,A}}{\rA}
  \proofstep{2}{M\subseteq\powerset(A)}{\rA}
  \proofstep{3}{X\in M}{\rA}
  \proofstep{1}{\QuotProj{A}{R}\colon A\to\QuotientSet{A}{R}}{%
    \FormulaRefAuto{QuotProjFunction}{1}}
  \proofstep{1}{\PowMap{\QuotProj{A}{R}}\colon
    \powerset(A)\to\powerset(\QuotientSet{A}{R})}{%
    \FormulaRefAuto{F\colon A\to B \vdash
    \PowMap{F}\colon \powerset(A)\to \powerset(B)}{4}}
  \proofstep{1,2}{\QuotFamMap{A}{R}{M}=
    \PowMap{\QuotProj{A}{R}}\restriction_M}{%
    \FormulaRefAuto{QuotFamMapDef}{1,2}}
  \proofstep{1,2,3}{\QuotFamMap{A}{R}{M}(X)=
    \bigl(\PowMap{\QuotProj{A}{R}}\restriction_M\bigr)(X)}{%
    \FormulaRefAuto{x\in A\dsep F=G\vdash F(x)=G(x)}{3,6}}
  \proofstep{1,2,3}{\bigl(\PowMap{\QuotProj{A}{R}}\restriction_M\bigr)(X)=
    \PowMap{\QuotProj{A}{R}}(X)}{%
    \FormulaRefAuto{C\subseteq A \dsep x\in C \vdash F\restriction_C(x)=F(x)}{2,3,5}}
  \proofstep{2,3}{X\in\powerset(A)}{%
    \FormulaRefAuto{A\subseteq B,\, x\in A \vdash x\in B}{2,3}}
  \proofstep{1,2,3}{\PowMap{\QuotProj{A}{R}}(X)=
    \QuotProj{A}{R}[X]}{%
    \FormulaRefAuto{F\colon A\to B \dsep X\in\powerset(A)\vdash
    \PowMap{F}(X)=F[X]}{4,9}}
  \proofstep{1,2,3}{\QuotFamMap{A}{R}{M}(X)=
    \QuotProj{A}{R}[X]}{%
    \FormulaRefAuto{a = b,\, b = c \vdash a = c}{7,
    \FormulaRefAuto{a = b,\, b = c \vdash a = c}{8,10}}}
\end{tabproof}

\subsubsection{Charakterisierung der Quotientenbildabbildung}

\FormulaThmDeltaK[Quotientenprojektionsbilder liegen in der Quotientenmenge]{%
  \EqRel{R,A}\dsep M\subseteq\powerset(A)\dsep X\in M
  \dsep C\in\QuotProj{A}{R}[X]
  \vdash
  C\in\QuotientSet{A}{R}
}{QuotProjImageInQuotient}{%
  \DeltaRow{Mengen}{A\dsep R\dsep M\dsep X\dsep C\dsep a}
  \DeltaRow{Funktionen}{\QuotProj{A}{R}}
  \DeltaRow{Quotientenmenge}{\QuotientSet{A}{R}}
}
\begin{tabproof}
  \proofstep{1}{\EqRel{R,A}}{\rA}
  \proofstep{2}{M\subseteq\powerset(A)}{\rA}
  \proofstep{3}{X\in M}{\rA}
  \proofstep{4}{C\in\QuotProj{A}{R}[X]}{\rA}
  \proofstep{2,3}{X\in\powerset(A)}{%
    \FormulaRefAuto{A\subseteq B,\, x\in A \vdash x\in B}{2,3}}
  \proofstep{2,3}{X\subseteq A}{%
    \FormulaRefAuto{x \in \powerset(A)\;\eqvdash\; x \subseteq A}{5}}
  \proofstep{1}{\QuotProj{A}{R}\colon
    A\to\QuotientSet{A}{R}}{%
    \FormulaRefAuto{QuotProjFunction}{1}}
  \proofstep{1,2,3}{\QuotProj{A}{R}[X]\subseteq
    \QuotientSet{A}{R}}{%
    \FormulaRefAuto{C\subseteq A \dsep F\colon A\to B \vdash F[C]\subseteq B}{6,7}}
  \proofstep{1,2,3,4}{C\in\QuotientSet{A}{R}}{%
    \FormulaRefAuto{A\subseteq B,\, x\in A \vdash x\in B}{8,4}}
\end{tabproof}

\FormulaThmDeltaK[Quotientenprojektionsbilder liefern Klassenzeugen]{%
  \EqRel{R,A}\dsep M\subseteq\powerset(A)\dsep X\in M
  \dsep C\in\QuotProj{A}{R}[X]
  \vdash
  \exists a\in X\,C=\EqClass{a}{R}{A}
}{QuotProjImageClassWitness}{%
  \DeltaRow{Mengen}{A\dsep R\dsep M\dsep X\dsep C\dsep a}
  \DeltaRow{Funktionen}{\QuotProj{A}{R}}
  \DeltaRow{Quotientenmenge}{\QuotientSet{A}{R}}
}
\begin{tabproofwide}
  \proofstepwidestar[1]{\EqRel{R,A}}{\rA}
  \proofstepwidestar[2]{M\subseteq\powerset(A)}{\rA}
  \proofstepwidestar[3]{X\in M}{\rA}
  \proofstepwidestar[4]{C\in\QuotProj{A}{R}[X]}{\rA}
  \proofstepwidestar[2,3]{X\in\powerset(A)}{%
    \FormulaRefAuto{A\subseteq B,\, x\in A \vdash x\in B}{2,3}}
  \proofstepwidestar[2,3]{X\subseteq A}{%
    \FormulaRefAuto{x \in \powerset(A)\;\eqvdash\; x \subseteq A}{5}}
  \proofstepwidestar[1]{\QuotProj{A}{R}\colon
    A\to\QuotientSet{A}{R}}{%
    \FormulaRefAuto{QuotProjFunction}{1}}
  \proofstepwidestar[1,2,3,4]{\exists a\in X\,
    C=\QuotProj{A}{R}(a)}{%
    \FormulaRefAuto{C\subseteq A\dsep F\colon A\to B\dsep
    y\in F[C]\vdash \exists x\in C\,y=F(x)}{6,7,4}}
  \proofstepwidestar[9]{a\in X\land C=\QuotProj{A}{R}(a)}{\rA}
  \proofstepwidestar[9]{a\in X}{\rAEa{9}}
  \proofstepwidestar[9]{C=\QuotProj{A}{R}(a)}{\rAEb{9}}
  \proofstepwidestar[2,3,9]{a\in A}{%
    \FormulaRefAuto{A\subseteq B,\, x\in A \vdash x\in B}{6,10}}
  \proofstepwidestar[1,9]{\QuotProj{A}{R}(a)=\EqClass{a}{R}{A}}{%
    \FormulaRefAuto{QuotProjValue}{1,12}}
  \proofstepwidestar[1,9]{C=\EqClass{a}{R}{A}}{%
    \FormulaRefAuto{a = b,\, b = c \vdash a = c}{11,13}}
  \proofstepwidestar[1,9]{a\in X\land C=\EqClass{a}{R}{A}}{\rAI{10,14}}
  \proofstepwidestar[1,9]{\exists a\in X\,C=\EqClass{a}{R}{A}}{\rEI{15}}
  \proofstepwidestar[1,2,3,4]{\exists a\in X\,C=\EqClass{a}{R}{A}}{%
    \rEE{8,9,16}}
\end{tabproofwide}

\FormulaThmDeltaK[Quotientenprojektionsbilder liefern Klassen]{%
  \EqRel{R,A}\dsep M\subseteq\powerset(A)\dsep X\in M
  \dsep C\in\QuotProj{A}{R}[X]
  \vdash
  C\in\QuotientSet{A}{R}\land
  \exists a\in X\,C=\EqClass{a}{R}{A}
}{QuotProjImageToEqClass}{%
  \DeltaRow{Mengen}{A\dsep R\dsep M\dsep X\dsep C\dsep a}
  \DeltaRow{Funktionen}{\QuotProj{A}{R}}
  \DeltaRow{Quotientenmenge}{\QuotientSet{A}{R}}
}
\begin{tabproof}
  \proofstep{1}{\EqRel{R,A}}{\rA}
  \proofstep{2}{M\subseteq\powerset(A)}{\rA}
  \proofstep{3}{X\in M}{\rA}
  \proofstep{4}{C\in\QuotProj{A}{R}[X]}{\rA}
  \proofstep{1,2,3,4}{C\in\QuotientSet{A}{R}}{%
    \FormulaRefAuto{QuotProjImageInQuotient}{1,2,3,4}}
  \proofstep{1,2,3,4}{\exists a\in X\,C=\EqClass{a}{R}{A}}{%
    \FormulaRefAuto{QuotProjImageClassWitness}{1,2,3,4}}
  \proofstep{1,2,3,4}{C\in\QuotientSet{A}{R}\land
    \exists a\in X\,C=\EqClass{a}{R}{A}}{\rAI{5,6}}
\end{tabproof}

\FormulaThmDeltaK[Klassen liefern Quotientenprojektionsbilder]{%
  \EqRel{R,A}\dsep M\subseteq\powerset(A)\dsep X\in M
  \dsep \exists a\in X\,C=\EqClass{a}{R}{A}
  \vdash
  C\in\QuotProj{A}{R}[X]
}{EqClassToQuotProjImage}{%
  \DeltaRow{Mengen}{A\dsep R\dsep M\dsep X\dsep C\dsep a}
  \DeltaRow{Funktionen}{\QuotProj{A}{R}}
  \DeltaRow{Quotientenmenge}{\QuotientSet{A}{R}}
}
\begin{tabproofwide}
  \proofstepwidestar[1]{\EqRel{R,A}}{\rA}
  \proofstepwidestar[2]{M\subseteq\powerset(A)}{\rA}
  \proofstepwidestar[3]{X\in M}{\rA}
  \proofstepwidestar[4]{\exists a\in X\,C=\EqClass{a}{R}{A}}{\rA}
  \proofstepwidestar[2,3]{X\in\powerset(A)}{%
    \FormulaRefAuto{A\subseteq B,\, x\in A \vdash x\in B}{2,3}}
  \proofstepwidestar[2,3]{X\subseteq A}{%
    \FormulaRefAuto{x \in \powerset(A)\;\eqvdash\; x \subseteq A}{5}}
  \proofstepwidestar[1]{\QuotProj{A}{R}\colon
    A\to\QuotientSet{A}{R}}{%
    \FormulaRefAuto{QuotProjFunction}{1}}
  \proofstepwidestar[8]{a\in X\land C=\EqClass{a}{R}{A}}{\rA}
  \proofstepwidestar[8]{a\in X}{\rAEa{8}}
  \proofstepwidestar[8]{C=\EqClass{a}{R}{A}}{\rAEb{8}}
  \proofstepwidestar[2,3,8]{a\in A}{%
    \FormulaRefAuto{A\subseteq B,\, x\in A \vdash x\in B}{6,9}}
  \proofstepwidestar[1,8]{\QuotProj{A}{R}(a)=\EqClass{a}{R}{A}}{%
    \FormulaRefAuto{QuotProjValue}{1,11}}
  \proofstepwidestar[1,8]{\EqClass{a}{R}{A}=\QuotProj{A}{R}(a)}{%
    \FormulaRefAuto{a=b\vdash b=a}{12}}
  \proofstepwidestar[1,8]{C=\QuotProj{A}{R}(a)}{%
    \FormulaRefAuto{a = b,\, b = c \vdash a = c}{10,13}}
  \proofstepwidestar[1,8]{a\in X\land C=\QuotProj{A}{R}(a)}{%
    \rAI{9,14}}
  \proofstepwidestar[1,8]{\exists a\in X\,
    C=\QuotProj{A}{R}(a)}{\rEI{15}}
  \proofstepwidestar[1,2,3,8]{C\in\QuotProj{A}{R}[X]}{%
    \FormulaRefAuto{C\subseteq A\dsep F\colon A\to B\dsep
    \exists x\in C\,y=F(x)\vdash y\in F[C]}{6,7,16}}
  \proofstepwidestar[1,2,3,4]{C\in\QuotProj{A}{R}[X]}{\rEE{4,8,17}}
\end{tabproofwide}

\FormulaThmDeltaK[Charakterisierung der Quotientenbildabbildung]{%
  \EqRel{R,A}\dsep M\subseteq\powerset(A)\dsep X\in M\vdash
  \begin{aligned}[t]
    C\in\QuotFamMap{A}{R}{M}(X)
    \leftrightarrow{}&
    C\in\QuotientSet{A}{R}\land{}\\
    &\exists a\in X\,C=\EqClass{a}{R}{A}
  \end{aligned}
}{QuotFamMapChar}{%
  \DeltaRow{Mengen}{A\dsep R\dsep M\dsep X\dsep C\dsep a}
  \DeltaRow{Funktionen}{\QuotProj{A}{R}\dsep \QuotFamMap{A}{R}{M}}
  \DeltaRow{Quotientenmenge}{\QuotientSet{A}{R}}
}
\begin{tabproofsplitwide}
  \proofpartwideindR[\(\rightarrow\)]{%
    \EqRel{R,A}\dsep M\subseteq\powerset(A)\dsep X\in M
    \vdash
    C\in\QuotFamMap{A}{R}{M}(X)\rightarrow
    \bigl(C\in\QuotientSet{A}{R}\land
    \exists a\in X\,C=\EqClass{a}{R}{A}\bigr)
  }{%
    \EqRel{R,A}\dsep M\subseteq\powerset(A)\dsep X\in M
    \vdash
    C\in\QuotFamMap{A}{R}{M}(X)\rightarrow
    \bigl(C\in\QuotientSet{A}{R}\land
    \exists a\in X\,C=\EqClass{a}{R}{A}\bigr)
  }[QuotFamMapCharForward]
    \proofstepwidestar[1]{\EqRel{R,A}}{\rA}
    \proofstepwidestar[2]{M\subseteq\powerset(A)}{\rA}
    \proofstepwidestar[3]{X\in M}{\rA}
    \proofstepwidestar[4]{C\in\QuotFamMap{A}{R}{M}(X)}{\rA}
    \proofstepwidestar[1,2,3]{\QuotFamMap{A}{R}{M}(X)=
      \QuotProj{A}{R}[X]}{%
      \FormulaRefAuto{QuotFamMapValue}{1,2,3}}
    \proofstepwidestar[1,2,3,4]{C\in\QuotProj{A}{R}[X]}{\rIE{5,4}}
    \proofstepwidestar[1,2,3,4]{C\in\QuotientSet{A}{R}\land
      \exists a\in X\,C=\EqClass{a}{R}{A}}{%
      \FormulaRefAuto{QuotProjImageToEqClass}{1,2,3,6}}
    \proofstepwidestar[1,2,3]{C\in\QuotFamMap{A}{R}{M}(X)
      \rightarrow
      \bigl(C\in\QuotientSet{A}{R}\land
      \exists a\in X\,C=\EqClass{a}{R}{A}\bigr)}{\rRI{4,7}}
  \closeproofpartwideindR

  \proofpartwideindR[\(\leftarrow\)]{%
    \EqRel{R,A}\dsep M\subseteq\powerset(A)\dsep X\in M
    \vdash
    \bigl(C\in\QuotientSet{A}{R}\land
    \exists a\in X\,C=\EqClass{a}{R}{A}\bigr)
    \rightarrow C\in\QuotFamMap{A}{R}{M}(X)
  }{%
    \EqRel{R,A}\dsep M\subseteq\powerset(A)\dsep X\in M
    \vdash
    \bigl(C\in\QuotientSet{A}{R}\land
    \exists a\in X\,C=\EqClass{a}{R}{A}\bigr)
    \rightarrow C\in\QuotFamMap{A}{R}{M}(X)
  }[QuotFamMapCharBackward]
    \proofstepwidestar[1]{\EqRel{R,A}}{\rA}
    \proofstepwidestar[2]{M\subseteq\powerset(A)}{\rA}
    \proofstepwidestar[3]{X\in M}{\rA}
    \proofstepwidestar[4]{C\in\QuotientSet{A}{R}\land
      \exists a\in X\,C=\EqClass{a}{R}{A}}{\rA}
    \proofstepwidestar[4]{\exists a\in X\,C=\EqClass{a}{R}{A}}{\rAEb{4}}
    \proofstepwidestar[1,2,3,4]{C\in\QuotProj{A}{R}[X]}{%
      \FormulaRefAuto{EqClassToQuotProjImage}{1,2,3,5}}
    \proofstepwidestar[1,2,3]{\QuotProj{A}{R}[X]=
      \QuotFamMap{A}{R}{M}(X)}{%
      \FormulaRefAuto{a=b\vdash b=a}{%
      \FormulaRefAuto{QuotFamMapValue}{1,2,3}}}
    \proofstepwidestar[1,2,3,4]{C\in\QuotFamMap{A}{R}{M}(X)}{%
      \rIE{7,6}}
    \proofstepwidestar[1,2,3]{%
      \bigl(C\in\QuotientSet{A}{R}\land
      \exists a\in X\,C=\EqClass{a}{R}{A}\bigr)
      \rightarrow C\in\QuotFamMap{A}{R}{M}(X)}{\rRI{4,8}}
  \closeproofpartwideindR

  \proofpartwide{Schluss}
    \proofstepwidestar[1]{\EqRel{R,A}}{\rA}
    \proofstepwidestar[2]{M\subseteq\powerset(A)}{\rA}
    \proofstepwidestar[3]{X\in M}{\rA}
    \proofstepwidestar[1,2,3]{C\in\QuotFamMap{A}{R}{M}(X)
      \rightarrow
      \bigl(C\in\QuotientSet{A}{R}\land
      \exists a\in X\,C=\EqClass{a}{R}{A}\bigr)}{%
      \FormulaRefAuto{QuotFamMapCharForward}{1,2,3}}
    \proofstepwidestar[1,2,3]{%
      \bigl(C\in\QuotientSet{A}{R}\land
      \exists a\in X\,C=\EqClass{a}{R}{A}\bigr)
      \rightarrow C\in\QuotFamMap{A}{R}{M}(X)}{%
      \FormulaRefAuto{QuotFamMapCharBackward}{1,2,3}}
    \proofstepwidestar[1,2,3]{C\in\QuotFamMap{A}{R}{M}(X)
      \leftrightarrow
      \bigl(C\in\QuotientSet{A}{R}\land
      \exists a\in X\,C=\EqClass{a}{R}{A}\bigr)}{\rLRI{4,5}}
  \closeproofpartwide
\end{tabproofsplitwide}

\FormulaThmDeltaK[Quotientenbildabbildungen erhalten Vereinigungen]{%
  \begin{aligned}[t]
    &\EqRel{R,A}\dsep M\subseteq\powerset(A)\dsep
      X\in M\dsep Y\in M\dsep X\cup Y\in M\\
    &\vdash
      \QuotFamMap{A}{R}{M}(X\cup Y)=
      \QuotFamMap{A}{R}{M}(X)\cup\QuotFamMap{A}{R}{M}(Y)
  \end{aligned}
}{QuotFamMapPreservesUnion}{%
  \DeltaRow{Mengen}{A\dsep R\dsep M\dsep X\dsep Y}
  \DeltaRow{Relationen}{R}
  \DeltaRow{Funktionen}{\QuotProj{A}{R}\dsep\QuotFamMap{A}{R}{M}}
}
\begin{tabproofwide}
  \proofstepwidestar[1]{\EqRel{R,A}}{\rA}
  \proofstepwidestar[2]{M\subseteq\powerset(A)}{\rA}
  \proofstepwidestar[3]{X\in M}{\rA}
  \proofstepwidestar[4]{Y\in M}{\rA}
  \proofstepwidestar[5]{X\cup Y\in M}{\rA}
  \proofstepwidestar[1]{\QuotProj{A}{R}\colon
    A\to\QuotientSet{A}{R}}{%
    \FormulaRefAuto{QuotProjFunction}{1}}
  \proofstepwidestar[2,3]{X\in\powerset(A)}{%
    \FormulaRefAuto{A\subseteq B,\,x\in A\vdash x\in B}{2,3}}
  \proofstepwidestar[2,3]{X\subseteq A}{%
    \FormulaRefAuto{x\in\powerset(A)\eqvdash x\subseteq A}{7}}
  \proofstepwidestar[2,4]{Y\in\powerset(A)}{%
    \FormulaRefAuto{A\subseteq B,\,x\in A\vdash x\in B}{2,4}}
  \proofstepwidestar[2,4]{Y\subseteq A}{%
    \FormulaRefAuto{x\in\powerset(A)\eqvdash x\subseteq A}{9}}
  \proofstepwidestar[1,2,3,4]{%
    \QuotProj{A}{R}[X\cup Y]=
    \QuotProj{A}{R}[X]\cup\QuotProj{A}{R}[Y]}{%
    \FormulaRefAuto{F\colon M\to N \dsep A\subseteq M \dsep
    B\subseteq M \vdash F[A\cup B]=F[A]\cup F[B]}{6,8,10}}
  \proofstepwidestar[1,2,5]{%
    \QuotFamMap{A}{R}{M}(X\cup Y)=\QuotProj{A}{R}[X\cup Y]}{%
    \FormulaRefAuto{QuotFamMapValue}{1,2,5}}
  \proofstepwidestar[1,2,3]{%
    \QuotFamMap{A}{R}{M}(X)=\QuotProj{A}{R}[X]}{%
    \FormulaRefAuto{QuotFamMapValue}{1,2,3}}
  \proofstepwidestar[1,2,4]{%
    \QuotFamMap{A}{R}{M}(Y)=\QuotProj{A}{R}[Y]}{%
    \FormulaRefAuto{QuotFamMapValue}{1,2,4}}
  \proofstepwidestar[1,2,3]{%
    \QuotProj{A}{R}[X]=\QuotFamMap{A}{R}{M}(X)}{%
    \FormulaRefAuto{a=b\vdash b=a}{13}}
  \proofstepwidestar[1,2,3]{%
    \QuotProj{A}{R}[X]\cup\QuotProj{A}{R}[Y]=
    \QuotFamMap{A}{R}{M}(X)\cup\QuotProj{A}{R}[Y]}{%
    \FormulaRefAuto{A=B\vdash A\cup C = B\cup C}{15}}
  \proofstepwidestar[1,2,4]{%
    \QuotProj{A}{R}[Y]=\QuotFamMap{A}{R}{M}(Y)}{%
    \FormulaRefAuto{a=b\vdash b=a}{14}}
  \proofstepwidestar[1,2,3,4]{%
    \QuotFamMap{A}{R}{M}(X)\cup\QuotProj{A}{R}[Y]=
    \QuotFamMap{A}{R}{M}(X)\cup\QuotFamMap{A}{R}{M}(Y)}{%
    \FormulaRefAuto{A=B\vdash C\cup A = C\cup B}{17}}
  \proofstepwidestar[1,2,3,4,5]{%
    \QuotFamMap{A}{R}{M}(X\cup Y)=
    \QuotProj{A}{R}[X]\cup\QuotProj{A}{R}[Y]}{%
    \FormulaRefAuto{a = b,\, b = c \vdash a = c}{12,11}}
  \proofstepwidestar[1,2,3,4,5]{%
    \QuotFamMap{A}{R}{M}(X\cup Y)=
    \QuotFamMap{A}{R}{M}(X)\cup\QuotProj{A}{R}[Y]}{%
    \FormulaRefAuto{a = b,\, b = c \vdash a = c}{19,16}}
  \proofstepwidestar[1,2,3,4,5]{%
    \QuotFamMap{A}{R}{M}(X\cup Y)=
    \QuotFamMap{A}{R}{M}(X)\cup\QuotFamMap{A}{R}{M}(Y)}{%
    \FormulaRefAuto{a = b,\, b = c \vdash a = c}{20,18}}
\end{tabproofwide}

\subsection{Ununterscheidbarkeitsquotient einer Mengenfamilie}

Die allgemeine Quotiententheorie wird nun auf die bereits eingeführte
Ununterscheidbarkeitsrelation einer Mengenfamilie spezialisiert. Die folgenden
Symbole sind nur Abkürzungen für die allgemeinen Konstruktionen.

\begin{notation*}
Für eine Mengenfamilie \(M\) schreiben wir
\[
  \begin{aligned}
    \OccClass{M}{a}
      &\coloneqq \EqClass{a}{\SameOccRel{M}}{\bigcup M},\\
    \OccQuotCarrier{M}
      &\coloneqq \QuotientSet{\bigcup M}{\SameOccRel{M}},\\
    \OccQuotMap{M}
      &\coloneqq \QuotFamMap{\bigcup M}{\SameOccRel{M}}{M},\\
    \OccQuotFam{M}
      &\coloneqq \OccQuotMap{M}[M].
  \end{aligned}
\]
\end{notation*}

Die folgenden kurzen Typisierungs- und Charakterisierungssätze dienen nur als
Schnittstelle für spätere Bände. Sie sind unmittelbare Spezialisierungen der
allgemeinen Sätze dieses Kapitels; neue inhaltliche Arbeit beginnt erst mit dem
Mitgliedschaftstransport.

\FormulaThmDeltaK[Typisierung der Ununterscheidbarkeitsquotientenabbildung]{%
  \OccQuotMap{M}\colon M\to\powerset(\OccQuotCarrier{M})
}{OccQuotMapPowerFunction}{%
  \DeltaRow{Mengen}{M}
  \DeltaRow{Relationen}{\SameOccRel{M}}
  \DeltaRow{Funktionen}{\OccQuotMap{M}}
}
\begin{tabproof}
  \proofstep{}{\EqRel{\SameOccRel{M},\bigcup M}}{%
    \FormulaRefAuto{SameOccRelEqRel}}
  \proofstep{}{M\subseteq\powerset(\bigcup M)}{%
    \FormulaRefAuto{FamilySubsetPowerUnion}}
  \proofstep{}{\OccQuotMap{M}\colon
    M\to\powerset(\OccQuotCarrier{M})}{%
    \FormulaRefAuto{QuotFamMapFunction}{1,2}}
\end{tabproof}

\FormulaThmDeltaK[Charakterisierung der Quotientenbilder]{%
  X\in M\vdash
  \begin{aligned}[t]
    C\in\OccQuotMap{M}(X)
    \leftrightarrow{}&
    C\in\OccQuotCarrier{M}\land{}\\
    &\exists a\in X\,C=\OccClass{M}{a}
  \end{aligned}
}{OccQuotMapChar}{%
  \DeltaRow{Mengen}{M\dsep X\dsep C\dsep a}
  \DeltaRow{Relationen}{\SameOccRel{M}}
  \DeltaRow{Funktionen}{\OccQuotMap{M}}
}
\begin{tabproof}
  \proofstep{1}{X\in M}{\rA}
  \proofstep{}{\EqRel{\SameOccRel{M},\bigcup M}}{%
    \FormulaRefAuto{SameOccRelEqRel}}
  \proofstep{}{M\subseteq\powerset(\bigcup M)}{%
    \FormulaRefAuto{FamilySubsetPowerUnion}}
  \proofstep{1}{%
    C\in\OccQuotMap{M}(X)\leftrightarrow
    \bigl(C\in\OccQuotCarrier{M}\land
    \exists a\in X\,C=\OccClass{M}{a}\bigr)}{%
    \FormulaRefAuto{QuotFamMapChar}{2,3,1}}
\end{tabproof}

\FormulaThmDeltaK[Typisierung auf die Quotientenfamilie]{%
  \OccQuotMap{M}\colon M\to\OccQuotFam{M}
}{OccQuotMapFunction}{%
  \DeltaRow{Mengen}{M}
  \DeltaRow{Funktionen}{\OccQuotMap{M}}
}
\begin{tabproof}
  \proofstep{}{\OccQuotMap{M}\colon
    M\to\powerset(\OccQuotCarrier{M})}{%
    \FormulaRefAuto{OccQuotMapPowerFunction}}
  \proofstep{}{\OccQuotMap{M}\colon M\sur\OccQuotFam{M}}{%
    \FormulaRefAuto{CorestrictionToImageSurjective}{1}}
  \proofstep{}{\OccQuotMap{M}\colon M\to\OccQuotFam{M}}{%
    \FormulaRefAuto{Surjektive Funktion}{2}}
\end{tabproof}

\FormulaThmDeltaK[Quotientenbilder gehören zur Quotientenfamilie]{%
  X\in M\vdash \OccQuotMap{M}(X)\in\OccQuotFam{M}
}{OccQuotSetInFam}{%
  \DeltaRow{Mengen}{M\dsep X}
  \DeltaRow{Funktionen}{\OccQuotMap{M}}
}
\begin{tabproof}
  \proofstep{1}{X\in M}{\rA}
  \proofstep{}{\OccQuotMap{M}\colon M\to\OccQuotFam{M}}{%
    \FormulaRefAuto{OccQuotMapFunction}}
  \proofstep{1}{\OccQuotMap{M}(X)\in\OccQuotFam{M}}{%
    \FormulaRefAuto{F\colon A\to B\dsep x\in A\vdash F(x)\in B}{2,1}}
\end{tabproof}

\FormulaThmDeltaK[Mitglieder der Quotientenfamilie besitzen Urbilder]{%
  Y\in\OccQuotFam{M}\vdash
  \exists X\in M\,Y=\OccQuotMap{M}(X)
}{OccQuotFamPreimage}{%
  \DeltaRow{Mengen}{M\dsep X\dsep Y}
  \DeltaRow{Funktionen}{\OccQuotMap{M}}
}
\begin{tabproof}
  \proofstep{1}{Y\in\OccQuotFam{M}}{\rA}
  \proofstep{}{\OccQuotMap{M}\colon
    M\to\powerset(\OccQuotCarrier{M})}{%
    \FormulaRefAuto{OccQuotMapPowerFunction}}
  \proofstep{1}{\exists X\in M\,Y=\OccQuotMap{M}(X)}{%
    \FormulaRefAuto{F\colon A\to B\dsep y\in F[A]\vdash
    \exists x\in A\,y=F(x)}{2,1}}
\end{tabproof}

\FormulaThmDeltaK[Mitglieder der Quotientenfamilie liegen im Träger]{%
  Y\in\OccQuotFam{M}\vdash Y\subseteq\OccQuotCarrier{M}
}{OccQuotMemberSubsetCarrier}{%
  \DeltaRow{Mengen}{M\dsep Y}
  \DeltaRow{Funktionen}{\OccQuotMap{M}}
}
\begin{tabproof}
  \proofstep{1}{Y\in\OccQuotFam{M}}{\rA}
  \proofstep{}{\OccQuotMap{M}\colon
    M\to\powerset(\OccQuotCarrier{M})}{%
    \FormulaRefAuto{OccQuotMapPowerFunction}}
  \proofstep{}{\OccQuotMap{M}[M]\subseteq
    \powerset(\OccQuotCarrier{M})}{%
    \FormulaRefAuto{F\colon A\to B\vdash F[A]\subseteq B}{2}}
  \proofstep{1}{Y\in\powerset(\OccQuotCarrier{M})}{%
    \FormulaRefAuto{A\subseteq B,\,x\in A\vdash x\in B}{3,1}}
  \proofstep{1}{Y\subseteq\OccQuotCarrier{M}}{%
    \FormulaRefAuto{x\in\powerset(A)\eqvdash x\subseteq A}{4}}
\end{tabproof}

\paragraph{Struktureigenschaften}

Ab hier werden die für Mengenfamilien wesentlichen Folgen der Spezialisierung
hergeleitet: Mitgliedschaftstransport, Trennung der Familienglieder,
Vereinigungsverträglichkeit und Bijektivität.

\FormulaThmDeltaK[Mitgliedschaftstransport in den Quotienten]{%
  X\in M\dsep a\in\bigcup M
  \vdash
  a\in X\leftrightarrow
  \OccClass{M}{a}\in\OccQuotMap{M}(X)
}{OccQuotMembershipTransport}{%
  \DeltaRow{Mengen}{M\dsep X\dsep a\dsep b}
  \DeltaRow{Relationen}{\SameOccRel{M}}
  \DeltaRow{Funktionen}{\OccQuotMap{M}}
}
\begin{tabproofwide}
  \proofstepwidestar[1]{X\in M}{\rA}
  \proofstepwidestar[2]{a\in\bigcup M}{\rA}
  \proofstepwidestar[]{\EqRel{\SameOccRel{M},\bigcup M}}{%
    \FormulaRefAuto{SameOccRelEqRel}}

  \proofstepwidestar[4]{a\in X}{\rA}
  \proofstepwidestar[2]{\OccClass{M}{a}\in\OccQuotCarrier{M}}{%
    \FormulaRefAuto{EqClassInQuotient}{3,2}}
  \proofstepwidestar[]{\OccClass{M}{a}=\OccClass{M}{a}}{\rII}
  \proofstepwidestar[4]{\exists b\in X\,
    \OccClass{M}{a}=\OccClass{M}{b}}{%
    \rEI{\rAI{4,6}}}
  \proofstepwidestar[2,4]{\OccClass{M}{a}\in\OccQuotCarrier{M}\land
    \exists b\in X\,\OccClass{M}{a}=\OccClass{M}{b}}{%
    \rAI{5,7}}
  \proofstepwidestar[1,2,4]{%
    \OccClass{M}{a}\in\OccQuotMap{M}(X)}{%
    \FormulaRefAuto{OccQuotMapChar}{1,8}}
  \proofstepwidestar[1,2]{%
    a\in X\rightarrow\OccClass{M}{a}\in\OccQuotMap{M}(X)}{%
    \rRI{4,9}}

  \proofstepwidestar[11]{%
    \OccClass{M}{a}\in\OccQuotMap{M}(X)}{\rA}
  \proofstepwidestar[1,11]{\OccClass{M}{a}\in\OccQuotCarrier{M}\land
    \exists b\in X\,\OccClass{M}{a}=\OccClass{M}{b}}{%
    \FormulaRefAuto{OccQuotMapChar}{1,11}}
  \proofstepwidestar[1,11]{\exists b\in X\,
    \OccClass{M}{a}=\OccClass{M}{b}}{\rAEb{12}}
  \proofstepwidestar[14]{b\in X\land
    \OccClass{M}{a}=\OccClass{M}{b}}{\rA}
  \proofstepwidestar[14]{b\in X}{\rAEa{14}}
  \proofstepwidestar[1,14]{b\in\bigcup M}{%
    \FormulaRefAuto{B\in A\dsep x\in B\vdash x\in\bigcup A}{1,15}}
  \proofstepwidestar[14]{\OccClass{M}{a}=\OccClass{M}{b}}{\rAEb{14}}
  \proofstepwidestar[1,2,14]{(a,b)\in\SameOccRel{M}}{%
    \FormulaRefAuto{EqClassRelFromEqual}{3,2,16,17}}
  \proofstepwidestar[1,2,14]{%
    a\in\bigcup M\land b\in\bigcup M\land
    \forall Z\in M\,(a\in Z\leftrightarrow b\in Z)}{%
    \FormulaRefAuto{SameOccRelChar}{18}}
  \proofstepwidestar[1,2,14]{%
    \forall Z\in M\,(a\in Z\leftrightarrow b\in Z)}{%
    \rAEb{\rAEb{19}}}
  \proofstepwidestar[1,2,14]{a\in X\leftrightarrow b\in X}{%
    \rRE{1,\rUE{20}}}
  \proofstepwidestar[1,2,14]{a\in X}{%
    \FormulaRefAuto{P\leftrightarrow Q,\,Q\vdash P}{21,15}}
  \proofstepwidestar[1,2,11]{a\in X}{%
    \rEE{13,14,22}}
  \proofstepwidestar[1,2]{%
    \OccClass{M}{a}\in\OccQuotMap{M}(X)\rightarrow a\in X}{%
    \rRI{11,23}}

  \proofstepwidestar[1,2]{%
    a\in X\leftrightarrow
    \OccClass{M}{a}\in\OccQuotMap{M}(X)}{%
    \rLRI{10,24}}
\end{tabproofwide}

\FormulaThmDeltaK[Quotientenbilder trennen Familienglieder]{%
  X\in M\dsep Y\in M\dsep
  \OccQuotMap{M}(X)=\OccQuotMap{M}(Y)
  \vdash X=Y
}{OccQuotSetSeparatesMembers}{%
  \DeltaRow{Mengen}{M\dsep X\dsep Y\dsep z}
  \DeltaRow{Funktionen}{\OccQuotMap{M}}
}
\begin{tabproofwide}
  \proofstepwidestar[1]{X\in M}{\rA}
  \proofstepwidestar[2]{Y\in M}{\rA}
  \proofstepwidestar[3]{\OccQuotMap{M}(X)=\OccQuotMap{M}(Y)}{\rA}

  \proofstepwidestar[4]{z\in X}{\rA}
  \proofstepwidestar[1,4]{z\in\bigcup M}{%
    \FormulaRefAuto{B\in A\dsep x\in B\vdash x\in\bigcup A}{1,4}}
  \proofstepwidestar[1,4]{%
    z\in X\leftrightarrow
    \OccClass{M}{z}\in\OccQuotMap{M}(X)}{%
    \FormulaRefAuto{OccQuotMembershipTransport}{1,5}}
  \proofstepwidestar[1,4]{%
    \OccClass{M}{z}\in\OccQuotMap{M}(X)}{%
    \FormulaRefAuto{P\leftrightarrow Q,\,P\vdash Q}{6,4}}
  \proofstepwidestar[1,3,4]{%
    \OccClass{M}{z}\in\OccQuotMap{M}(Y)}{\rIE{3,7}}
  \proofstepwidestar[1,2,4]{%
    z\in Y\leftrightarrow
    \OccClass{M}{z}\in\OccQuotMap{M}(Y)}{%
    \FormulaRefAuto{OccQuotMembershipTransport}{2,5}}
  \proofstepwidestar[1,2,3,4]{z\in Y}{%
    \FormulaRefAuto{P\leftrightarrow Q,\,Q\vdash P}{9,8}}
  \proofstepwidestar[1,2,3]{X\subseteq Y}{%
    \FormulaRefAuto{A\subseteq B\coloneq
    \forall z\,(z\in A\rightarrow z\in B)}{\rUI{\rRI{4,10}}}}

  \proofstepwidestar[12]{z\in Y}{\rA}
  \proofstepwidestar[2,12]{z\in\bigcup M}{%
    \FormulaRefAuto{B\in A\dsep x\in B\vdash x\in\bigcup A}{2,12}}
  \proofstepwidestar[2,12]{%
    z\in Y\leftrightarrow
    \OccClass{M}{z}\in\OccQuotMap{M}(Y)}{%
    \FormulaRefAuto{OccQuotMembershipTransport}{2,13}}
  \proofstepwidestar[2,12]{%
    \OccClass{M}{z}\in\OccQuotMap{M}(Y)}{%
    \FormulaRefAuto{P\leftrightarrow Q,\,P\vdash Q}{14,12}}
  \proofstepwidestar[3]{\OccQuotMap{M}(Y)=\OccQuotMap{M}(X)}{%
    \FormulaRefAuto{a=b\vdash b=a}{3}}
  \proofstepwidestar[2,3,12]{%
    \OccClass{M}{z}\in\OccQuotMap{M}(X)}{\rIE{16,15}}
  \proofstepwidestar[1,2,12]{%
    z\in X\leftrightarrow
    \OccClass{M}{z}\in\OccQuotMap{M}(X)}{%
    \FormulaRefAuto{OccQuotMembershipTransport}{1,13}}
  \proofstepwidestar[1,2,3,12]{z\in X}{%
    \FormulaRefAuto{P\leftrightarrow Q,\,Q\vdash P}{18,17}}
  \proofstepwidestar[1,2,3]{Y\subseteq X}{%
    \FormulaRefAuto{A\subseteq B\coloneq
    \forall z\,(z\in A\rightarrow z\in B)}{\rUI{\rRI{12,19}}}}
  \proofstepwidestar[1,2,3]{X=Y}{%
    \FormulaRefAuto{A\subseteq B, B\subseteq A\vdash A=B}{11,20}}
\end{tabproofwide}

\FormulaThmDeltaK[Vereinigungsverträglichkeit des Ununterscheidbarkeitsquotienten]{%
  X\in M\dsep Y\in M\dsep X\cup Y\in M
  \vdash
  \OccQuotMap{M}(X\cup Y)=
  \OccQuotMap{M}(X)\cup\OccQuotMap{M}(Y)
}{OccQuotPreservesUnion}{%
  \DeltaRow{Mengen}{M\dsep X\dsep Y}
  \DeltaRow{Funktionen}{\OccQuotMap{M}}
}
\begin{tabproof}
  \proofstep{1}{X\in M}{\rA}
  \proofstep{2}{Y\in M}{\rA}
  \proofstep{3}{X\cup Y\in M}{\rA}
  \proofstep{}{\EqRel{\SameOccRel{M},\bigcup M}}{%
    \FormulaRefAuto{SameOccRelEqRel}}
  \proofstep{}{M\subseteq\powerset(\bigcup M)}{%
    \FormulaRefAuto{FamilySubsetPowerUnion}}
  \proofstep{1,2,3}{%
    \OccQuotMap{M}(X\cup Y)=
    \OccQuotMap{M}(X)\cup\OccQuotMap{M}(Y)}{%
    \FormulaRefAuto{QuotFamMapPreservesUnion}{4,5,1,2,3}}
\end{tabproof}

\FormulaThmDeltaK[Bijektivität der Ununterscheidbarkeitsquotientenabbildung]{%
  \OccQuotMap{M}\colon M\bij\OccQuotFam{M}
}{OccQuotMapBijective}{%
  \DeltaRow{Mengen}{M\dsep X\dsep Y}
  \DeltaRow{Funktionen}{\OccQuotMap{M}}
}
\begin{tabproof}
  \proofstep{}{\OccQuotMap{M}\colon M\to\OccQuotFam{M}}{%
    \FormulaRefAuto{OccQuotMapFunction}}
  \proofstep{2}{X\in M}{\rA}
  \proofstep{3}{Y\in M}{\rA}
  \proofstep{4}{\OccQuotMap{M}(X)=\OccQuotMap{M}(Y)}{\rA}
  \proofstep{2,3,4}{X=Y}{%
    \FormulaRefAuto{OccQuotSetSeparatesMembers}{2,3,4}}
  \proofstep{}{\OccQuotMap{M}\colon M\inj\OccQuotFam{M}}{%
    \FormulaRefAuto{Injektive Funktion}{1,5}}
  \proofstep{}{\OccQuotMap{M}\colon
    M\to\powerset(\OccQuotCarrier{M})}{%
    \FormulaRefAuto{OccQuotMapPowerFunction}}
  \proofstep{}{\OccQuotMap{M}\colon M\sur\OccQuotFam{M}}{%
    \FormulaRefAuto{CorestrictionToImageSurjective}{7}}
  \proofstep{}{\OccQuotMap{M}\colon M\bij\OccQuotFam{M}}{%
    \FormulaRefAuto{F\colon A\inj B\dsep F\colon A\sur B
    \vdash F\colon A\bij B}{6,8}}
\end{tabproof}

\FormulaThmDeltaK[Vereinigungsabschluss der Quotientenfamilie]{%
  \UnionClosedFam(M)\vdash\UnionClosedFam(\OccQuotFam{M})
}{OccQuotUnionClosed}{%
  \DeltaRow{Mengen}{M\dsep X\dsep Y\dsep U\dsep V}
  \DeltaRow{Funktionen}{\OccQuotMap{M}}
}
\begin{tabproofwide}
  \proofstepwidestar[1]{\UnionClosedFam(M)}{\rA}
  \proofstepwidestar[2]{U\in\OccQuotFam{M}}{\rA}
  \proofstepwidestar[3]{V\in\OccQuotFam{M}}{\rA}
  \proofstepwidestar[2]{\exists X\in M\,U=\OccQuotMap{M}(X)}{%
    \FormulaRefAuto{OccQuotFamPreimage}{2}}
  \proofstepwidestar[3]{\exists Y\in M\,V=\OccQuotMap{M}(Y)}{%
    \FormulaRefAuto{OccQuotFamPreimage}{3}}
  \proofstepwidestar[6]{X\in M\land U=\OccQuotMap{M}(X)}{\rA}
  \proofstepwidestar[7]{Y\in M\land V=\OccQuotMap{M}(Y)}{\rA}
  \proofstepwidestar[6]{X\in M}{\rAEa{6}}
  \proofstepwidestar[7]{Y\in M}{\rAEa{7}}
  \proofstepwidestar[1,6,7]{X\cup Y\in M}{%
    \FormulaRefAuto{Vereinigungsabgeschlossene Familie}{1,8,9}}
  \proofstepwidestar[6,7]{%
    U\cup V=\OccQuotMap{M}(X)\cup\OccQuotMap{M}(Y)}{%
    \rIE{\rAEb{6},\rAEb{7}}}
  \proofstepwidestar[1,6,7]{%
    \OccQuotMap{M}(X)\cup\OccQuotMap{M}(Y)=
    \OccQuotMap{M}(X\cup Y)}{%
    \FormulaRefAuto{a=b\vdash b=a}{%
      \FormulaRefAuto{OccQuotPreservesUnion}{8,9,10}}}
  \proofstepwidestar[1,6,7]{%
    U\cup V=\OccQuotMap{M}(X\cup Y)}{%
    \FormulaRefAuto{a = b,\, b = c \vdash a = c}{11,12}}
  \proofstepwidestar[1,6,7]{%
    \OccQuotMap{M}(X\cup Y)\in\OccQuotFam{M}}{%
    \FormulaRefAuto{OccQuotSetInFam}{10}}
  \proofstepwidestar[1,6,7]{U\cup V\in\OccQuotFam{M}}{%
    \rIE{13,14}}
  \proofstepwidestar[1,6]{U\cup V\in\OccQuotFam{M}}{%
    \rEE{5,7,15}}
  \proofstepwidestar[1]{U\cup V\in\OccQuotFam{M}}{%
    \rEE{4,6,16}}
  \proofstepwidestar[1]{\UnionClosedFam(\OccQuotFam{M})}{%
    \FormulaRefAuto{Vereinigungsabgeschlossene Familie}{%
      \rUI{\rRI{2,\rRI{3,17}}}}}
\end{tabproofwide}

\FormulaThmDeltaK[Vereinigung der Quotientenfamilie]{%
  \bigcup\OccQuotFam{M}=\OccQuotCarrier{M}
}{OccQuotFamUnionCarrier}{%
  \DeltaRow{Mengen}{M\dsep X\dsep Y\dsep C\dsep a}
  \DeltaRow{Relationen}{\SameOccRel{M}}
  \DeltaRow{Funktionen}{\OccQuotMap{M}}
}
\begin{tabproofwide}
  \proofstepwidestar[1]{C\in\bigcup\OccQuotFam{M}}{\rA}
  \proofstepwidestar[1]{\exists Y\in\OccQuotFam{M}\,C\in Y}{%
    \FormulaRefAuto{x\in\bigcup A\eqvdash
    \exists B\in A\,(x\in B)}{1}}
  \proofstepwidestar[3]{Y\in\OccQuotFam{M}\land C\in Y}{\rA}
  \proofstepwidestar[3]{Y\in\OccQuotFam{M}}{\rAEa{3}}
  \proofstepwidestar[3]{C\in Y}{\rAEb{3}}
  \proofstepwidestar[3]{\exists X\in M\,Y=\OccQuotMap{M}(X)}{%
    \FormulaRefAuto{OccQuotFamPreimage}{4}}
  \proofstepwidestar[7]{X\in M\land Y=\OccQuotMap{M}(X)}{\rA}
  \proofstepwidestar[7]{X\in M}{\rAEa{7}}
  \proofstepwidestar[3,7]{C\in\OccQuotMap{M}(X)}{%
    \rIE{\rAEb{7},5}}
  \proofstepwidestar[3,7]{C\in\OccQuotCarrier{M}\land
    \exists a\in X\,C=\OccClass{M}{a}}{%
    \FormulaRefAuto{OccQuotMapChar}{8,9}}
  \proofstepwidestar[3,7]{C\in\OccQuotCarrier{M}}{\rAEa{10}}
  \proofstepwidestar[3]{C\in\OccQuotCarrier{M}}{%
    \rEE{6,7,11}}
  \proofstepwidestar[1]{C\in\OccQuotCarrier{M}}{%
    \rEE{2,3,12}}
  \proofstepwidestar[]{%
    \bigcup\OccQuotFam{M}\subseteq\OccQuotCarrier{M}}{%
    \FormulaRefAuto{A\subseteq B\coloneq
    \forall C\,(C\in A\rightarrow C\in B)}{\rUI{\rRI{1,13}}}}

  \proofstepwidestar[15]{C\in\OccQuotCarrier{M}}{\rA}
  \proofstepwidestar[]{\EqRel{\SameOccRel{M},\bigcup M}}{%
    \FormulaRefAuto{SameOccRelEqRel}}
  \proofstepwidestar[15]{\exists a\in\bigcup M\,
    C=\OccClass{M}{a}}{%
    \FormulaRefAuto{QuotientRepresentative}{16,15}}
  \proofstepwidestar[18]{a\in\bigcup M\land C=\OccClass{M}{a}}{\rA}
  \proofstepwidestar[18]{a\in\bigcup M}{\rAEa{18}}
  \proofstepwidestar[18]{C=\OccClass{M}{a}}{\rAEb{18}}
  \proofstepwidestar[18]{\exists X\in M\,a\in X}{%
    \FormulaRefAuto{x\in\bigcup A\eqvdash
    \exists B\in A\,(x\in B)}{19}}
  \proofstepwidestar[22]{X\in M\land a\in X}{\rA}
  \proofstepwidestar[22]{X\in M}{\rAEa{22}}
  \proofstepwidestar[22]{a\in X}{\rAEb{22}}
  \proofstepwidestar[18,22]{%
    a\in X\leftrightarrow
    \OccClass{M}{a}\in\OccQuotMap{M}(X)}{%
    \FormulaRefAuto{OccQuotMembershipTransport}{23,19}}
  \proofstepwidestar[18,22]{%
    \OccClass{M}{a}\in\OccQuotMap{M}(X)}{%
    \FormulaRefAuto{P\leftrightarrow Q,\,P\vdash Q}{25,24}}
  \proofstepwidestar[18,22]{C\in\OccQuotMap{M}(X)}{%
    \rIE{20,26}}
  \proofstepwidestar[22]{\OccQuotMap{M}(X)\in\OccQuotFam{M}}{%
    \FormulaRefAuto{OccQuotSetInFam}{23}}
  \proofstepwidestar[18,22]{C\in\bigcup\OccQuotFam{M}}{%
    \FormulaRefAuto{x\in\bigcup A\eqvdash
    \exists B\in A\,(x\in B)}{\rEI{\rAI{28,27}}}}
  \proofstepwidestar[18]{C\in\bigcup\OccQuotFam{M}}{%
    \rEE{21,22,29}}
  \proofstepwidestar[15]{C\in\bigcup\OccQuotFam{M}}{%
    \rEE{17,18,30}}
  \proofstepwidestar[]{%
    \OccQuotCarrier{M}\subseteq\bigcup\OccQuotFam{M}}{%
    \FormulaRefAuto{A\subseteq B\coloneq
    \forall C\,(C\in A\rightarrow C\in B)}{\rUI{\rRI{15,31}}}}
  \proofstepwidestar[]{%
    \bigcup\OccQuotFam{M}=\OccQuotCarrier{M}}{%
    \FormulaRefAuto{A\subseteq B, B\subseteq A\vdash A=B}{14,32}}
\end{tabproofwide}

\FormulaThmDeltaK[Elemente der vereinigten Quotientenfamilie haben Repräsentanten]{%
  C\in\bigcup\OccQuotFam{M}
  \vdash
  \exists a\in\bigcup M\,C=\OccClass{M}{a}
}{OccQuotUnionRepresentative}{%
  \DeltaRow{Mengen}{M\dsep C\dsep a}
  \DeltaRow{Relationen}{\SameOccRel{M}}
}
\begin{tabproof}
  \proofstep{1}{C\in\bigcup\OccQuotFam{M}}{\rA}
  \proofstep{}{\bigcup\OccQuotFam{M}=\OccQuotCarrier{M}}{%
    \FormulaRefAuto{OccQuotFamUnionCarrier}}
  \proofstep{1}{C\in\OccQuotCarrier{M}}{\rIE{2,1}}
  \proofstep{}{\EqRel{\SameOccRel{M},\bigcup M}}{%
    \FormulaRefAuto{SameOccRelEqRel}}
  \proofstep{1}{\exists a\in\bigcup M\,C=\OccClass{M}{a}}{%
    \FormulaRefAuto{QuotientRepresentative}{4,3}}
\end{tabproof}

\FormulaThmDeltaK[Nichtleerheit der Quotientenfamilie]{%
  M\neq\varnothing\vdash\OccQuotFam{M}\neq\varnothing
}{OccQuotFamNonempty}{%
  \DeltaRow{Mengen}{M\dsep X}
  \DeltaRow{Funktionen}{\OccQuotMap{M}}
}
\begin{tabproof}
  \proofstep{1}{M\neq\varnothing}{\rA}
  \proofstep{1}{\exists X\in M}{%
    \FormulaRefAuto{B\neq\varnothing\vdash\exists x\in B}{1}}
  \proofstep{3}{X\in M}{\rA}
  \proofstep{3}{\OccQuotMap{M}(X)\in\OccQuotFam{M}}{%
    \FormulaRefAuto{OccQuotSetInFam}{3}}
  \proofstep{3}{\OccQuotFam{M}\neq\varnothing}{%
    \FormulaRefAuto{a\in A\vdash A\neq\varnothing}{4}}
  \proofstep{1}{\OccQuotFam{M}\neq\varnothing}{\rEE{2,3,5}}
\end{tabproof}

\FormulaThmDeltaK[Nichtleerheit der Vereinigung der Quotientenfamilie]{%
  \bigcup M\neq\varnothing
  \vdash
  \bigcup\OccQuotFam{M}\neq\varnothing
}{OccQuotFamCarrierNonempty}{%
  \DeltaRow{Mengen}{M\dsep a}
  \DeltaRow{Relationen}{\SameOccRel{M}}
}
\begin{tabproof}
  \proofstep{1}{\bigcup M\neq\varnothing}{\rA}
  \proofstep{1}{\exists a\in\bigcup M}{%
    \FormulaRefAuto{B\neq\varnothing\vdash\exists x\in B}{1}}
  \proofstep{3}{a\in\bigcup M}{\rA}
  \proofstep{}{\EqRel{\SameOccRel{M},\bigcup M}}{%
    \FormulaRefAuto{SameOccRelEqRel}}
  \proofstep{3}{\OccClass{M}{a}\in\OccQuotCarrier{M}}{%
    \FormulaRefAuto{EqClassInQuotient}{4,3}}
  \proofstep{}{\bigcup\OccQuotFam{M}=\OccQuotCarrier{M}}{%
    \FormulaRefAuto{OccQuotFamUnionCarrier}}
  \proofstep{}{\OccQuotCarrier{M}=\bigcup\OccQuotFam{M}}{%
    \FormulaRefAuto{a=b\vdash b=a}{6}}
  \proofstep{3}{\OccClass{M}{a}\in\bigcup\OccQuotFam{M}}{%
    \rIE{7,5}}
  \proofstep{3}{\bigcup\OccQuotFam{M}\neq\varnothing}{%
    \FormulaRefAuto{a\in A\vdash A\neq\varnothing}{8}}
  \proofstep{1}{\bigcup\OccQuotFam{M}\neq\varnothing}{%
    \rEE{2,3,9}}
\end{tabproof}

\paragraph{Die Spur trennt Klassen}

Die allgemeine Spurabbildung zeichnet jede Quotientenklasse durch die
Familienglieder aus, in denen sie vorkommt. Für den
Ununterscheidbarkeitsquotienten ist diese Kennzeichnung injektiv.

\FormulaThmDeltaK[Die allgemeine Spur trennt Ununterscheidbarkeitsklassen]{%
  \TraceMap{\OccQuotMap{M}}\colon
  \OccQuotCarrier{M}\inj\powerset(M)
}{OccQuotTraceMapInjective}{%
  \DeltaRow{Mengen}{M\dsep C\dsep D\dsep X\dsep a\dsep b}
  \DeltaRow{Relationen}{\SameOccRel{M}}
  \DeltaRow{Funktionen}{\OccQuotMap{M}\dsep
    \TraceMap{\OccQuotMap{M}}}
}
\begin{tabproofwide}
  \proofstepwidestar[]{\OccQuotMap{M}\colon
    M\to\powerset(\OccQuotCarrier{M})}{%
    \FormulaRefAuto{OccQuotMapPowerFunction}}
  \proofstepwidestar[]{\TraceMap{\OccQuotMap{M}}\colon
    \OccQuotCarrier{M}\to\powerset(M)}{%
    \FormulaRefAuto{TraceMapFunction}{1}}
  \proofstepwidestar[3]{C\in\OccQuotCarrier{M}}{\rA}
  \proofstepwidestar[4]{D\in\OccQuotCarrier{M}}{\rA}
  \proofstepwidestar[5]{%
    \TraceMap{\OccQuotMap{M}}(C)=
    \TraceMap{\OccQuotMap{M}}(D)}{\rA}
  \proofstepwidestar[]{\EqRel{\SameOccRel{M},\bigcup M}}{%
    \FormulaRefAuto{SameOccRelEqRel}}
  \proofstepwidestar[3]{\exists a\in\bigcup M\,
    C=\OccClass{M}{a}}{%
    \FormulaRefAuto{QuotientRepresentative}{6,3}}
  \proofstepwidestar[4]{\exists b\in\bigcup M\,
    D=\OccClass{M}{b}}{%
    \FormulaRefAuto{QuotientRepresentative}{6,4}}

  \proofstepwidestar[9]{a\in\bigcup M\land
    C=\OccClass{M}{a}}{\rA}
  \proofstepwidestar[9]{a\in\bigcup M}{\rAEa{9}}
  \proofstepwidestar[9]{C=\OccClass{M}{a}}{\rAEb{9}}
  \proofstepwidestar[12]{b\in\bigcup M\land
    D=\OccClass{M}{b}}{\rA}
  \proofstepwidestar[12]{b\in\bigcup M}{\rAEa{12}}
  \proofstepwidestar[12]{D=\OccClass{M}{b}}{\rAEb{12}}

  \proofstepwidestar[15]{X\in M}{\rA}
  \proofstepwidestar[3,4,5,15]{%
    C\in\OccQuotMap{M}(X)\leftrightarrow
    D\in\OccQuotMap{M}(X)}{%
    \FormulaRefAuto{TraceMapEqualityMembership}{1,3,4,5,15}}
  \proofstepwidestar[9,15]{%
    a\in X\leftrightarrow
    \OccClass{M}{a}\in\OccQuotMap{M}(X)}{%
    \FormulaRefAuto{OccQuotMembershipTransport}{15,10}}
  \proofstepwidestar[9,15]{%
    a\in X\leftrightarrow C\in\OccQuotMap{M}(X)}{%
    \rIE{11,17}}
  \proofstepwidestar[12,15]{%
    b\in X\leftrightarrow
    \OccClass{M}{b}\in\OccQuotMap{M}(X)}{%
    \FormulaRefAuto{OccQuotMembershipTransport}{15,13}}
  \proofstepwidestar[12,15]{%
    b\in X\leftrightarrow D\in\OccQuotMap{M}(X)}{%
    \rIE{14,19}}
  \proofstepwidestar[12,15]{%
    D\in\OccQuotMap{M}(X)\leftrightarrow b\in X}{%
    \FormulaRefAuto{P\leftrightarrow Q\vdash Q\leftrightarrow P}{20}}
  \proofstepwidestar[3,4,5,9,15]{%
    a\in X\leftrightarrow D\in\OccQuotMap{M}(X)}{%
    \FormulaRefAuto{P\leftrightarrow Q\dsep Q\leftrightarrow R
    \vdash P\leftrightarrow R}{18,16}}
  \proofstepwidestar[3,4,5,9,12,15]{a\in X\leftrightarrow b\in X}{%
    \FormulaRefAuto{P\leftrightarrow Q\dsep Q\leftrightarrow R
    \vdash P\leftrightarrow R}{22,21}}
  \proofstepwidestar[3,4,5,9,12]{%
    \forall X\in M\,(a\in X\leftrightarrow b\in X)}{%
    \rUI{\rRI{15,23}}}
  \proofstepwidestar[3,4,5,9,12]{%
    a\in\bigcup M\land b\in\bigcup M\land
    \forall X\in M\,(a\in X\leftrightarrow b\in X)}{%
    \rAI{10,\rAI{13,24}}}
  \proofstepwidestar[3,4,5,9,12]{(a,b)\in\SameOccRel{M}}{%
    \FormulaRefAuto{SameOccRelChar}{25}}
  \proofstepwidestar[3,4,5,9,12]{%
    \OccClass{M}{a}=\OccClass{M}{b}}{%
    \FormulaRefAuto{EqClassEqualFromRel}{6,10,13,26}}
  \proofstepwidestar[3,4,5,9,12]{C=\OccClass{M}{b}}{%
    \FormulaRefAuto{a = b,\, b = c \vdash a = c}{11,27}}
  \proofstepwidestar[3,4,5,9,12]{C=D}{%
    \FormulaRefAuto{a = b,\, c = b \vdash a = c}{28,14}}
  \proofstepwidestar[3,4,5,9]{C=D}{\rEE{8,12,29}}
  \proofstepwidestar[3,4,5]{C=D}{\rEE{7,9,30}}
  \proofstepwidestar[]{\TraceMap{\OccQuotMap{M}}\colon
    \OccQuotCarrier{M}\inj\powerset(M)}{%
    \FormulaRefAuto{Injektive Funktion}{2,31}}
\end{tabproofwide}

\paragraph{Beispiel.}
Für \(M=\{\varnothing,\{a,b\}\}\) treten \(a\) und \(b\) in genau denselben
Familiengliedern auf. Daher bilden sie dieselbe Klasse
\(C=\OccClass{M}{a}=\OccClass{M}{b}\). Es gilt
\[
  \OccQuotMap{M}(\varnothing)=\varnothing,\qquad
  \OccQuotMap{M}(\{a,b\})=\{C\},
  \qquad
  \OccQuotFam{M}=\{\varnothing,\{C\}\}.
\]


\section{Gleichmächtigkeit}

\FormulaDefDeltaK[Begriff der gleichmächtigen Mengen]{A \EqCard B \coloneqq \exists F\colon A\bij B}{Gleichmächtigkeit}{
  \DeltaRow{Mengen}{A\dsep B}
  \DeltaRow{\textbf{Neue Symbole}}{}
  \DeltaRow{Gleichmächtigkeit}{}
}

\FormulaThmDeltaK[Frische Adjunktion erhält Gleichmächtigkeit]{%
  A\EqCard B\dsep a\notin A\dsep b\notin B
  \vdash A\cup\{a\}\EqCard B\cup\{b\}%
}{EqCardFreshAdjunction}{%
  \DeltaRow{Mengen}{A\dsep B\dsep a\dsep b}%
  \DeltaRow{Funktionen}{F}%
}
\begin{tabproof}
  \proofstep{1}{A\EqCard B}{\rA}
  \proofstep{2}{a\notin A}{\rA}
  \proofstep{3}{b\notin B}{\rA}
  \proofstep{1}{\exists F\colon A\bij B}{%
    \FormulaRefAuto{A \EqCard B \coloneqq \exists F\colon A\bij B}{1}}
  \proofstep{5}{F\colon A\bij B}{\rA}
  \proofstep{2,3,5}{%
    \ExtMap{\iota_{B,B\cup\{b\}}\circ F}{a}{b}
    \colon A\cup\{a\}\bij B\cup\{b\}}{%
    \FormulaRefAuto{FreshExtMapBijective}[thm]{5,2,3}}
  \proofstep{2,3,5}{\exists G\colon A\cup\{a\}\bij B\cup\{b\}}{\rEI{6}}
  \proofstep{2,3,5}{A\cup\{a\}\EqCard B\cup\{b\}}{%
    \FormulaRefAuto{A \EqCard B \coloneqq \exists F\colon A\bij B}{7}}
  \proofstep{1,2,3}{A\cup\{a\}\EqCard B\cup\{b\}}{\rEE{4,5,8}}
\end{tabproof}

\FormulaThmDeltaK[Reflexivität der Gleichmächtigkeit]{%
  A \EqCard A%
}{EqCardReflexive}{
  \DeltaRow{Mengen}{A}
}
\begin{tabproof}
  \proofstep{}{\Id_A\colon A\bij A}{\FormulaRefAuto{\Id_A\colon A\bij A}}
  \proofstep{}{\exists F\colon A\bij A}{\rEI{1}}
  \proofstep{}{A \EqCard A}{\FormulaRefAuto{A \EqCard B \coloneqq \exists F\colon A\bij B}{2}}
\end{tabproof}

\FormulaThmDeltaK[Symmetrie der Gleichmächtigkeit]{%
  A \EqCard B \vdash B \EqCard A%
}{EqCardSymmetric}{
  \DeltaRow{Mengen}{A\dsep B}
  \DeltaRow{Funktionen}{F\dsep G}
}
\begin{tabproof}
  \proofstep{1}{A \EqCard B}{\rA}
  \proofstep{1}{\exists F\colon A\bij B}{\FormulaRefAuto{A \EqCard B \coloneqq \exists F\colon A\bij B}{1}}

  \proofstep{3}{F\colon A\bij B}{\rA}
  \proofstep{3}{F^{-1}\colon B\bij A}{\FormulaRefAuto{F\colon A\bij B\vdash F^{-1}\colon B\bij A}{3}}
  \proofstep{3}{\exists G\colon B\bij A}{\rEI{4}}

  \proofstep{1}{\exists G\colon B\bij A}{\rEE{2,3,5}}
  \proofstep{1}{B \EqCard A}{\FormulaRefAuto{A \EqCard B \coloneqq \exists F\colon A\bij B}{6}}
\end{tabproof}

\FormulaThmDeltaK[Transitivität der Gleichmächtigkeit]{%
  A \EqCard B \dsep B \EqCard C \vdash A \EqCard C%
}{EqCardTransitive}{
  \DeltaRow{Mengen}{A\dsep B\dsep C}
  \DeltaRow{Funktionen}{F\dsep G\dsep H}
}
\begin{tabproof}
  \proofstep{1}{A \EqCard B}{\rA}
  \proofstep{2}{B \EqCard C}{\rA}

  \proofstep{1}{\exists F\colon A\bij B}{\FormulaRefAuto{A \EqCard B \coloneqq \exists F\colon A\bij B}{1}}
  \proofstep{2}{\exists G\colon B\bij C}{\FormulaRefAuto{A \EqCard B \coloneqq \exists F\colon A\bij B}{2}}

  \proofstep{5}{F\colon A\bij B}{\rA}
  \proofstep{6}{G\colon B\bij C}{\rA}
  \proofstep{5,6}{G\circ F\colon A\bij C}{\FormulaRefAuto{F\colon A\bij B \dsep G\colon B\bij C\vdash G\circ F\colon A\bij C}{5,6}}
  \proofstep{5,6}{\exists H\colon A\bij C}{\rEI{7}}
  \proofstep{5}{\exists H\colon A\bij C}{\rEE{4,6,8}}
  \proofstep{1,2}{\exists H\colon A\bij C}{\rEE{3,5,9}}

  \proofstep{1,2}{A \EqCard C}{\FormulaRefAuto{A \EqCard B \coloneqq \exists F\colon A\bij B}{10}}
\end{tabproof}

\FormulaThmDeltaK[Gleichmächtigkeit ist eine Äquivalenzrelation]{%
  \begin{aligned}[t]
    \vdash{}& A\EqCard A\land{}\\
    &\bigl(A\EqCard B\rightarrow B\EqCard A\bigr)\land{}\\
    &\bigl((A\EqCard B\land B\EqCard C)\rightarrow A\EqCard C\bigr)
  \end{aligned}
}{EqCardEquivalence}{%
  \DeltaRow{Mengen}{A\dsep B\dsep C}
}
\begin{tabproofwide}
  \proofstepwidestar[]{A\EqCard A}{%
    \FormulaRefAuto{EqCardReflexive}}
  \proofstepwidestar[2]{A\EqCard B}{\rA}
  \proofstepwidestar[2]{B\EqCard A}{%
    \FormulaRefAuto{EqCardSymmetric}{2}}
  \proofstepwidestar[]{A\EqCard B\rightarrow B\EqCard A}{%
    \rRI{2,3}}
  \proofstepwidestar[5]{A\EqCard B\land B\EqCard C}{\rA}
  \proofstepwidestar[5]{A\EqCard B}{\rAEa{5}}
  \proofstepwidestar[5]{B\EqCard C}{\rAEb{5}}
  \proofstepwidestar[5]{A\EqCard C}{%
    \FormulaRefAuto{EqCardTransitive}{6,7}}
  \proofstepwidestar[]{%
    (A\EqCard B\land B\EqCard C)\rightarrow A\EqCard C}{%
    \rRI{5,8}}
  \proofstepwidestar[]{%
    A\EqCard A\land
    \bigl(A\EqCard B\rightarrow B\EqCard A\bigr)\land
    \bigl((A\EqCard B\land B\EqCard C)\rightarrow A\EqCard C\bigr)}{%
    \rAI{1,\rAI{4,9}}}
\end{tabproofwide}

\FormulaThmDeltaK[Transport von Injektionen entlang der Gleichmächtigkeit]{%
  A \EqCard M\dsep B \EqCard N\dsep F\colon A\inj B
  \vdash \exists G\colon M\inj N
}{EqCardInjectionTransport}{%
  \DeltaRow{Mengen}{A\dsep B\dsep M\dsep N}
  \DeltaRow{Funktionen}{F\dsep G\dsep H\dsep K}
}
\begin{tabproof}
  \proofstep{1}{A \EqCard M}{\rA}
  \proofstep{2}{B \EqCard N}{\rA}
  \proofstep{3}{F\colon A\inj B}{\rA}
  \proofstep{1}{\exists H\colon A\bij M}{%
    \FormulaRefAuto{A \EqCard B \coloneqq \exists F\colon A\bij B}{1}}
  \proofstep{2}{\exists K\colon B\bij N}{%
    \FormulaRefAuto{A \EqCard B \coloneqq \exists F\colon A\bij B}{2}}

  \proofstep{6}{H\colon A\bij M}{\rA}
  \proofstep{7}{K\colon B\bij N}{\rA}
  \proofstep{6}{H^{-1}\colon M\bij A}{%
    \FormulaRefAuto{F\colon A\bij B\vdash F^{-1}\colon B\bij A}{6}}
  \proofstep{6}{H^{-1}\colon M\inj A}{%
    \FormulaRefAuto{F\colon A\bij B \vdash F\colon A\inj B}{8}}
  \proofstep{3,6}{F\circ H^{-1}\colon M\inj B}{%
    \FormulaRefAuto{F\colon A\inj B \dsep G\colon B\inj C\vdash G\circ F\colon A\inj C}{9,3}}
  \proofstep{7}{K\colon B\inj N}{%
    \FormulaRefAuto{F\colon A\bij B \vdash F\colon A\inj B}{7}}
  \proofstep{3,6,7}{K\circ(F\circ H^{-1})\colon M\inj N}{%
    \FormulaRefAuto{F\colon A\inj B \dsep G\colon B\inj C\vdash G\circ F\colon A\inj C}{10,11}}
  \proofstep{3,6,7}{\exists G\colon M\inj N}{\rEI{12}}

  \proofstep{2,3,6}{\exists G\colon M\inj N}{\rEE{5,7,13}}
  \proofstep{1,2,3}{\exists G\colon M\inj N}{\rEE{4,6,14}}
\end{tabproof}


\FormulaThmDelta[Potenzmengen gleichmächtiger Mengen sind gleichmächtig]{%
  A \EqCard B \vdash \powerset(A)\EqCard \powerset(B)
}{%
  \DeltaRow{Mengen}{A\dsep B}
  \DeltaRow{Funktionen}{F\dsep G\dsep \widehat{F}}
}
\begin{tabproof}
  \proofstep{1}{A \EqCard B}{\rA}
  \proofstep{1}{\exists F\colon A\bij B}{%
    \FormulaRefAuto{A \EqCard B \coloneqq \exists F\colon A\bij B}{1}}

  \proofstep{3}{F\colon A\bij B}{\rA}
    \proofstep{3}{\PowMap{F}\colon \powerset(A)\bij \powerset(B)}{%
      \FormulaRefAuto{F\colon A\bij B \vdash \PowMap{F}\colon \powerset(A)\bij \powerset(B)}{3}}
  \proofstep{3}{\exists G\colon \powerset(A)\bij \powerset(B)}{\rEI{4}}

  \proofstep{1}{\exists G\colon \powerset(A)\bij \powerset(B)}{\rEE{2,3,5}}
  \proofstep{1}{\powerset(A)\EqCard \powerset(B)}{%
    \FormulaRefAuto{A \EqCard B \coloneqq \exists F\colon A\bij B}{6}}
\end{tabproof}


\end{document}
